\documentclass[11pt,a4paper]{article}

% ============================================================
% PACOTES
% ============================================================

\usepackage[T1]{fontenc}
\usepackage[utf8]{inputenc}
\usepackage[brazil]{babel}

\usepackage{lmodern}
\usepackage{microtype}

\usepackage{amsmath,amssymb,amsthm,mathtools}
\usepackage{enumitem}

\usepackage{geometry}
\usepackage{setspace}

\usepackage{xcolor}
\usepackage{hyperref}
\usepackage{cleveref}

% ============================================================
% FORMATAÇÃO
% ============================================================

\geometry{
    top=2.7cm,
    bottom=2.7cm,
    left=2.7cm,
    right=2.7cm
}

\onehalfspacing

\setlength{\parindent}{1.25cm}
\setlength{\parskip}{0pt}

\hypersetup{
    colorlinks=true,
    linkcolor=blue!50!black,
    citecolor=blue!50!black,
    urlcolor=blue!50!black
}

% ============================================================
% AMBIENTES
% ============================================================

\newtheorem{theorem}{Teorema}[section]
\newtheorem{proposition}[theorem]{Proposição}
\newtheorem{lemma}[theorem]{Lema}
\newtheorem{corollary}[theorem]{Corolário}
\newtheorem{conjecture}[theorem]{Conjectura}

\theoremstyle{definition}
\newtheorem{definition}[theorem]{Definição}
\newtheorem{question}[theorem]{Questão}
\newtheorem{problem}[theorem]{Problema}

\theoremstyle{remark}
\newtheorem{remark}[theorem]{Observação}

% ============================================================
% COMANDOS
% ============================================================

\newcommand{\Gnp}{G(n,p)}
\newcommand{\whp}{\text{a.a.s.}}
\newcommand{\N}{\mathbb{N}}
\newcommand{\R}{\mathbb{R}}

\newcommand{\dens}{d}
\newcommand{\maxdens}{m}
\newcommand{\dtwo}{d_2}
\newcommand{\mtwo}{m_2}

\newcommand{\can}{\mathrm{can}}
\newcommand{\rb}{\mathrm{rb}}
\newcommand{\ram}{\mathrm{Ram}}

% ============================================================
% DOCUMENTO
% ============================================================

\begin{document}

% ============================================================
% CABEÇALHO
% ============================================================

\begin{center}

  {\Large\bfseries PROJETO DE DOUTORADO}

  \vspace{0.7cm}

  {\LARGE\bfseries
    LIMIARES PARA PROPRIEDADES RAMSEY CANÔNICAS\\
    E ANTI-RAMSEY EM GRAFOS ALEATÓRIOS
  }

  \vspace{1.2cm}

  \begin{tabular}{rl}
    \textbf{Candidato:}
     & Afonso Lima dos Santos Sant'Anna             \\[0.15cm]

    \textbf{Orientador:}
     & Yoshiharu Kohayakawa                         \\[0.15cm]

    \textbf{Coorientador:}
     & Guilherme Oliveira Mota                      \\[0.15cm]

    \textbf{Instituição:}
     & Instituto de Matemática e Estatística -- USP
  \end{tabular}

\end{center}

\vspace{0.8cm}

% ============================================================
% RESUMO
% ============================================================

\noindent
\textbf{Resumo.}
Este projeto de doutorado tem como objetivo investigar limiares para
propriedades do tipo Ramsey em grafos aleatórios, com ênfase em problemas
nos quais as colorações das arestas não estão sujeitas à restrição clássica
de utilizar um número fixo de cores.
O projeto será desenvolvido em três linhas principais.
A primeira consiste no estudo de propriedades anti-Ramsey, buscando
caracterizar classes de grafos para as quais o limiar é governado pela
$2$-densidade máxima.
A segunda trata de uma versão assimétrica da propriedade Ramsey canônica,
na qual os grafos procurados nas configurações monocromática, arco-íris e
lexicográfica podem ser distintos.
A terceira linha investiga limiares Ramsey canônicos para grafos além de
cliques, com particular interesse em ciclos e em condições estruturais que
garantam que o limiar canônico coincida com a escala Ramsey clássica.
Apesar de distintas, essas linhas são fortemente relacionadas por uma
mesma estratégia: reduzir problemas probabilísticos sobre $\Gnp$ a
problemas determinísticos de coloração em grafos de densidade controlada.
Pretendemos utilizar técnicas provenientes do método probabilístico, da
teoria de grafos aleatórios, de decomposições de grafos, degenerescência,
arboricidade, métodos de containers e estruturas de blocos associadas a
cópias dos grafos-alvo.

\vspace{0.4cm}

\tableofcontents

\newpage

% ============================================================
% INTRODUÇÃO
% ============================================================

\section{Introdução}

A Teoria de Ramsey estuda o aparecimento inevitável de estruturas
organizadas em sistemas suficientemente grandes. Em sua formulação clássica
para grafos, fixamos um grafo $H$ e um número finito de cores e perguntamos
quão grande deve ser um grafo $G$ para que toda coloração de suas arestas
contenha uma cópia monocromática de $H$. A partir dessa questão surgiram
diversas variantes, nas quais se modificam tanto a estrutura ambiente quanto
as restrições impostas à coloração.

Uma direção particularmente fértil consiste em substituir o grafo ambiente
determinístico por um grafo aleatório. Denotamos por $\Gnp$ o grafo
binomial em $n$ vértices, no qual cada possível aresta é escolhida
independentemente com probabilidade $p=p(n)$. Para uma propriedade
monótona crescente $\mathcal{P}$, uma questão fundamental é determinar a
escala de $p$ na qual $\Gnp$ passa de não satisfazer $\mathcal{P}$,
assintoticamente quase certamente, para satisfazê-la.

No contexto Ramsey clássico, esse problema foi essencialmente resolvido por
Rödl e Ruciński. Um dos aspectos centrais dessa teoria é o papel
desempenhado pela \emph{$2$-densidade máxima}
\[
  \mtwo(H)
  :=
  \max_{\substack{J\subseteq H\\ v(J)\geq 3}}
  \frac{e(J)-1}{v(J)-2}.
\]
Para uma ampla classe de grafos $H$, o limiar para que toda coloração das
arestas de $\Gnp$ com um número fixo de cores contenha uma cópia
monocromática de $H$ encontra-se na escala
\[
  p=n^{-1/\mtwo(H)}.
\]

Esse fenômeno sugere uma conexão profunda entre propriedades determinísticas
de coloração, parâmetros de densidade e a estrutura local das cópias de $H$
presentes no grafo aleatório. Em particular, os $0$-enunciados de resultados
Ramsey aleatórios frequentemente possuem uma contraparte determinística:
grafos cuja densidade máxima é suficientemente pequena admitem colorações
que evitam as configurações Ramsey desejadas.

Essa filosofia foi desenvolvida de maneira sistemática em trabalhos
posteriores. Em particular, o arcabouço algorítmico de Nenadov, Person,
Škorić e Steger permite, em diversas situações, reduzir um problema
probabilístico sobre $\Gnp$ à existência de determinadas obstruções
determinísticas de tamanho e densidade limitados. Em linhas gerais, o grafo
aleatório é submetido a um procedimento de remoção de estruturas facilmente
coloríveis; o núcleo remanescente é decomposto em blocos de tamanho
controlado; e a parte probabilística do argumento é então reduzida a um
problema determinístico de coloração desses blocos.

O presente projeto pretende investigar o que ocorre quando abandonamos
algumas das restrições da Teoria de Ramsey clássica. A primeira direção
considerada corresponde à propriedade \emph{anti-Ramsey}. Nesse contexto,
consideramos colorações próprias das arestas e procuramos cópias
\emph{arco-íris} de um grafo $H$, isto é, cópias cujas arestas recebem cores
duas a duas distintas. Uma questão natural é determinar para quais funções
$p=p(n)$ toda coloração própria das arestas de $\Gnp$ contém,
assintoticamente quase certamente, uma cópia arco-íris de $H$.

Em diversos casos, o limiar anti-Ramsey é novamente governado por
$\mtwo(H)$. Isso ocorre, por exemplo, para cliques de ordem pelo menos $5$ e
para ciclos de comprimento pelo menos $5$. Entretanto, o fenômeno não é
universal. Grafos pequenos, como $K_4$ e $C_4$, apresentam limiares
distintos da escala $n^{-1/\mtwo(H)}$. Essas exceções indicam que, ao
contrário do problema Ramsey clássico, ainda não existe uma caracterização
estrutural satisfatória dos grafos cujo limiar anti-Ramsey é determinado
pela $2$-densidade máxima.

Uma das linhas centrais deste projeto será, portanto, compreender quais
características estruturais do grafo $H$ garantem que
\[
  p_H^{\rb}(n)
  =
  n^{-1/\mtwo(H)}
\]
em ordem de grandeza. Em particular, investigaremos inicialmente grafos de
cintura grande, para os quais a ausência de ciclos curtos pode fornecer
decomposições mais favoráveis e permitir a construção de colorações próprias
sem cópias arco-íris do grafo-alvo.

Uma segunda direção surge quando não se impõe qualquer restrição ao número
de cores utilizado. Nesse cenário, uma coloração pode evitar
trivialmente cópias monocromáticas atribuindo, por exemplo, uma cor
distinta a cada aresta. A Teoria Canônica de Ramsey, iniciada por Erdős e
Rado, mostra que a ausência de uma restrição no número de cores não elimina
toda forma de estrutura: em cliques suficientemente grandes aparecem
necessariamente padrões especiais de coloração, entre eles cópias
monocromáticas, arco-íris e lexicográficas.

O estudo de colorações canônicas já esteve no centro do trabalho de
mestrado do candidato, no qual foram investigadas classes Ramsey em
hipergrafos e construções canônicas por meio do método partido. O presente
projeto dá continuidade natural a essa direção, agora com ênfase em
problemas probabilísticos.

Recentemente, Kamčev e Schacht determinaram o limiar Ramsey canônico para
cliques no grafo aleatório. Para todo $\ell\geq4$, o limiar para que toda
coloração arbitrária das arestas de $\Gnp$ contenha uma cópia canônica de
$K_\ell$ encontra-se na escala
\[
  n^{-1/\mtwo(K_\ell)}
  =
  n^{-2/(\ell+1)}.
\]
Assim, para cliques de ordem pelo menos $4$, permitir um número arbitrário
de cores não altera a escala do limiar Ramsey clássico.

Esse resultado levanta uma questão consideravelmente mais ampla:
\emph{para quais grafos $H$ o limiar Ramsey canônico é governado por
  $\mtwo(H)$?} Mesmo para classes fundamentais de grafos, como ciclos, o
cenário ainda não está completamente compreendido. Uma das linhas deste
projeto será investigar limiares Ramsey canônicos para grafos além de
cliques, com ênfase inicial em ciclos e em classes de grafos para as quais
as técnicas probabilísticas e estruturais disponíveis possam ser
adaptadas.

Há ainda uma terceira possibilidade, situada entre o problema canônico
simétrico e as diferentes variantes assimétricas da Teoria de Ramsey.
Dados três grafos $H_1,H_2,H_3$, podemos perguntar pelo limiar para que
toda coloração arbitrária das arestas de $\Gnp$ contenha pelo menos uma
das seguintes configurações:
\begin{enumerate}[label=\textup{(\roman*)}]
  \item uma cópia monocromática de $H_1$;
  \item uma cópia arco-íris de $H_2$;
  \item uma cópia lexicográfica de $H_3$.
\end{enumerate}
Chamaremos essa propriedade de \emph{Ramsey canônica assimétrica}.

Essa formulação é inspirada tanto no teorema canônico de Erdős--Rado quanto
nos problemas Ramsey assimétricos, nos quais diferentes cores possuem
diferentes grafos-alvo. No problema Ramsey assimétrico clássico, a
interação entre grafos de densidades distintas leva naturalmente à
introdução de parâmetros de densidade mistos, como aquele que aparece na
conjectura de Kohayakawa--Kreuter. Além disso, avanços recentes mostram
que grande parte desses $0$-enunciados pode ser reduzida a problemas
determinísticos de coloração e decomposição.

Esses resultados sugerem procurar um fenômeno análogo no contexto
canônico: determinar quais parâmetros governam a interação entre as
alternativas monocromática, arco-íris e lexicográfica e investigar até que
ponto os respectivos problemas probabilísticos podem ser reduzidos a
condições determinísticas de densidade.

Um primeiro objetivo nessa direção será estudar sistematicamente triplas
de cliques
\[
  (K_r,K_s,K_t).
\]
Os casos de cliques fornecem um ambiente suficientemente estruturado para
que diferentes técnicas de coloração possam ser testadas e, ao mesmo
tempo, apresentam fenômenos não triviais quando as ordens dos três alvos
são distintas. A partir desses primeiros casos, pretendemos investigar
formulações mais gerais envolvendo outros grafos e buscar parâmetros de
densidade adequados à versão canônica assimétrica.

As três direções descritas acima compartilham, portanto, uma mesma
filosofia. Em cada uma delas, pretende-se compreender como
\[
  \boxed{
    \text{estrutura do grafo-alvo}
    \;+\;
    \text{restrições sobre a coloração}
    \;+\;
    \text{densidade}
  }
\]
determinam o limiar de uma propriedade do tipo Ramsey em $\Gnp$.

Em particular, um tema recorrente será a transformação de $0$-enunciados
probabilísticos em problemas determinísticos sobre grafos de densidade
limitada. Para isso, pretendemos explorar ferramentas como degenerescência,
arboricidade e decomposições em florestas, estruturas de núcleos e blocos,
argumentos de primeiro e segundo momento, métodos probabilísticos, técnicas
de containers e construções de coloração adaptadas aos padrões que se
deseja evitar.

O projeto será organizado em três linhas principais de investigação:
\begin{enumerate}
  \item \textbf{Limiar anti-Ramsey:}
        investigar condições estruturais sob as quais
        \[
          p_H^{\rb}(n)=n^{-1/\mtwo(H)},
        \]
        com ênfase inicial em grafos de cintura grande;

  \item \textbf{Ramsey canônico assimétrico:}
        estudar inicialmente triplas de cliques $(K_r,K_s,K_t)$ e,
        posteriormente, buscar parâmetros de densidade e técnicas que permitam
        tratar grafos mais gerais;

  \item \textbf{Ramsey canônico para grafos gerais:}
        investigar para quais grafos $H$ o limiar canônico coincide com
        $n^{-1/\mtwo(H)}$, tendo ciclos como uma das primeiras classes de
        interesse.
\end{enumerate}

Essas linhas contêm problemas de diferentes níveis de dificuldade. Algumas
questões iniciais reduzem-se a construções explícitas de colorações em
grafos de densidade limitada; outras exigem adaptações de frameworks
probabilísticos conhecidos; e os objetivos mais ambiciosos buscam
identificar princípios estruturais capazes de explicar os limiares
observados em classes amplas de grafos.

% ============================================================
% PRÓXIMAS SEÇÕES
% ============================================================

% ============================================================
% OBJETIVOS E JUSTIFICATIVA
% ============================================================

\section{Objetivos e justificativa}

O objetivo geral deste projeto é investigar limiares para propriedades
do tipo Ramsey em grafos aleatórios, com particular interesse em
situações nas quais as colorações das arestas não estão restritas a um
número fixo de cores. Pretendemos compreender de que maneira parâmetros
de densidade do grafo-alvo, propriedades estruturais das colorações e
interações entre diferentes configurações Ramsey determinam o
comportamento assintótico dessas propriedades em $G(n,p)$.

Um fenômeno recorrente na teoria Ramsey aleatória é o aparecimento da
$2$-densidade máxima
\[
  m_2(H)
  =
  \max_{\substack{J\subseteq H\\v(J)\geq 3}}
  \frac{e(J)-1}{v(J)-2}
\]
como parâmetro determinante do limiar. No problema Ramsey clássico,
esse fenômeno é amplamente compreendido: para uma extensa classe de
grafos $H$, o limiar encontra-se na escala
\[
  n^{-1/m_2(H)}.
\]
Em variantes mais recentes, entretanto, o panorama é consideravelmente
menos uniforme. Tanto no problema anti-Ramsey quanto no problema
Ramsey canônico existem resultados que apontam novamente para a
$2$-densidade como parâmetro natural, mas também aparecem exceções,
problemas ainda em aberto e situações em que a interação entre
diferentes grafos-alvo sugere a necessidade de novos parâmetros.

Uma das motivações centrais deste projeto é compreender até que ponto
a filosofia desenvolvida para o problema Ramsey clássico pode ser
estendida a essas variantes. Em particular, pretendemos explorar
situações nas quais um problema probabilístico sobre $G(n,p)$ possa ser
reduzido a um problema determinístico de coloração em grafos de
densidade limitada. Esse princípio está presente, por exemplo, no
framework algorítmico de Nenadov, Person, Škorić e Steger, no qual
$0$-enunciados para propriedades Ramsey em grafos aleatórios são
reduzidos à análise de obstruções determinísticas de tamanho e
densidade controlados. Esse tipo de redução será uma das ideias
norteadoras do projeto.

Nesse contexto, o projeto será desenvolvido em três direções
principais, que embora distintas são fortemente relacionadas.

\subsection*{Objetivo I: limiares anti-Ramsey}

O primeiro objetivo é avançar na compreensão do limiar para a
propriedade anti-Ramsey em grafos aleatórios. Dados grafos $G$ e $H$,
consideramos a propriedade de que toda coloração própria das arestas de
$G$ contenha uma cópia arco-íris de $H$.

Para diversas classes importantes de grafos, como cliques e ciclos
suficientemente grandes, o limiar conhecido está na escala
\[
  n^{-1/m_2(H)}.
\]
Entretanto, grafos pequenos como $K_4$ e $C_4$ apresentam comportamento
distinto, mostrando que a $2$-densidade, isoladamente, não fornece uma
caracterização geral.

Nosso objetivo será identificar condições estruturais sobre $H$ que
garantam que
\[
  p_H^{\mathrm{rb}}(n)
  =
  n^{-1/m_2(H)}
\]
em ordem de grandeza. Uma primeira classe de interesse será formada por
grafos de cintura grande. A ausência de ciclos curtos pode restringir
fortemente as possíveis interações entre diferentes cópias de $H$ e
permitir a construção de colorações próprias em grafos de densidade
limitada.

Mais amplamente, pretendemos compreender quais estruturas locais
produzem as exceções conhecidas ao comportamento governado por
$m_2(H)$ e buscar critérios determinísticos que distingam grafos
``regulares'' do ponto de vista anti-Ramsey daqueles que apresentam
obstruções de baixa densidade.

\subsection*{Objetivo II: Ramsey canônico assimétrico}

O segundo objetivo é desenvolver uma teoria probabilística para uma
versão assimétrica do problema Ramsey canônico.

Dados três grafos ordenados $H_1,H_2,H_3$, pretendemos estudar a
propriedade
\[
  G(n,p)
  \xrightarrow{\mathrm{can}}
  (H_1,H_2,H_3),
\]
segundo a qual toda coloração arbitrária das arestas de $G(n,p)$ contém
uma cópia monocromática de $H_1$, uma cópia arco-íris de $H_2$ ou uma
cópia lexicográfica de $H_3$.

O caso em que
\[
  H_1=H_2=H_3
\]
corresponde à propriedade Ramsey canônica usual. Quando os três grafos
são distintos, entretanto, surge um fenômeno genuinamente assimétrico:
cada uma das três alternativas pode ser governada por parâmetros de
densidade diferentes, e a interação entre essas escalas pode produzir
um comportamento mais complexo.

Pretendemos iniciar esse estudo pela família
\[
  (K_r,K_s,K_t),
\]
para a qual a estrutura dos cliques permite testar diferentes
estratégias de coloração e formular conjecturas precisas para os
limiares. Resultados preliminares sugerem que, em diversos regimes, a
ordem do maior clique exerce papel dominante, mas também indicam a
existência de casos pequenos excepcionais, em particular quando
$K_3$ ou $K_4$ aparecem entre os alvos.

Um objetivo de maior alcance será identificar um parâmetro de densidade
adequado à versão canônica assimétrica, possivelmente análogo aos
parâmetros mistos que aparecem na teoria Ramsey assimétrica clássica.
Também investigaremos até que ponto os $0$-enunciados desse problema
podem ser reduzidos a questões determinísticas sobre grafos de
densidade limitada.

\subsection*{Objetivo III: limiares Ramsey canônicos para grafos gerais}

O terceiro objetivo é investigar a propriedade Ramsey canônica em
grafos aleatórios para grafos-alvo que não sejam necessariamente
completos.

O resultado de Kamčev e Schacht para cliques mostra que, para todo
$\ell\geq4$, o limiar Ramsey canônico de $K_\ell$ está na escala
\[
  n^{-1/m_2(K_\ell)}.
\]
Esse resultado sugere perguntar em que medida o mesmo fenômeno ocorre
para outras classes de grafos.

Pretendemos estudar inicialmente ciclos, que constituem uma das classes
mais naturais para testar essa questão. Em particular, interessa
determinar o limiar canônico para ciclos ímpares, compreender com maior
precisão o caso dos ciclos pares e investigar se a escala
\[
  n^{-1/m_2(H)}
\]
continua sendo a escala correta.

Em uma perspectiva mais geral, pretendemos buscar condições estruturais
sobre $H$ que garantam que o limiar Ramsey canônico coincida com o
limiar Ramsey clássico. Entre as propriedades que poderão ser
investigadas estão $2$-balanceamento, cintura, bipartição,
degenerescência e outras condições de esparsidade ou regularidade.

\subsection*{Justificativa e integração das linhas de pesquisa}

As três linhas propostas não constituem problemas independentes.
Todas elas podem ser vistas como manifestações de uma mesma questão
fundamental: determinar quais propriedades estruturais de um grafo e
de uma classe de colorações controlam o aparecimento inevitável de
configurações Ramsey em grafos aleatórios.

Um aspecto particularmente importante do projeto será explorar a
relação entre argumentos probabilísticos e problemas determinísticos
de coloração. O framework de Nenadov, Person, Škorić e Steger mostra
que, em diversas propriedades Ramsey, o estudo do $0$-enunciado pode
ser reduzido à seguinte questão: dado um grafo $B$ cuja densidade
máxima está abaixo de um determinado parâmetro crítico, é possível
colorir suas arestas evitando a configuração proibida?
Esse paradigma está diretamente relacionado às três linhas deste
projeto.

No problema anti-Ramsey, isso leva à busca por colorações próprias
sem cópias arco-íris. No problema Ramsey canônico assimétrico, conduz
à construção de colorações que evitem simultaneamente padrões
monocromáticos, arco-íris e lexicográficos. Finalmente, no problema
Ramsey canônico usual, pode fornecer uma ferramenta para estender a
teoria atualmente conhecida para classes de grafos mais gerais.

As técnicas necessárias para essas questões incluem métodos
probabilísticos, argumentos de primeiro e segundo momento,
decomposições de grafos, degenerescência, arboricidade, orientações
com grau de saída controlado, estruturas de núcleos e blocos e
técnicas de containers. Em várias situações, problemas probabilísticos
serão acompanhados por lemas puramente determinísticos de coloração,
permitindo uma interação constante entre combinatória extremal e
probabilidade.

Além dos problemas centrais, o projeto foi concebido de forma a
oferecer diferentes níveis de dificuldade e múltiplas possibilidades
de desenvolvimento. Algumas questões admitem inicialmente o estudo de
casos particulares, como cliques pequenos ou ciclos de comprimento
fixo, enquanto outras permitem formular conjecturas estruturais mais
amplas. Essa organização possibilita que avanços parciais produzam
resultados independentes e, ao mesmo tempo, contribuam para os
objetivos gerais do projeto.

Por fim, o projeto representa uma continuação natural da formação do
candidato em Teoria de Ramsey e Combinatória Extremal. O trabalho de
mestrado foi dedicado ao estudo de colorações canônicas em hipergrafos
e ao uso de construções do tipo partido. O presente projeto amplia
essa perspectiva ao contexto probabilístico e busca conectar técnicas
de Ramsey canônico, anti-Ramsey e grafos aleatórios em um único
programa de pesquisa.

% ============================================================
% TÓPICOS ESPECÍFICOS DE PESQUISA
% ============================================================

\section{Tópicos específicos de pesquisa}

\subsection{Limiar anti-Ramsey em grafos aleatórios}
\label{sec:anti-ramsey}

A primeira linha de pesquisa deste projeto trata de propriedades
anti-Ramsey em grafos aleatórios. Dados grafos $G$ e $H$, escrevemos
\[
  G\xrightarrow{\mathrm{rb}} H
\]
quando toda coloração própria das arestas de $G$ contém uma cópia
arco-íris de $H$. Denotaremos por
\[
  p_H^{\mathrm{rb}}=p_H^{\mathrm{rb}}(n)
\]
o limiar correspondente no grafo aleatório $G(n,p)$.

O estudo sistemático dessa propriedade no contexto aleatório foi iniciado
por Kohayakawa, Konstadinidis e Mota, que provaram que, para todo grafo
fixo $H$, existe uma constante $C>0$ tal que
\[
  p\geq Cn^{-1/m_2(H)}
\]
implica
\[
  G(n,p)\xrightarrow{\mathrm{rb}}H
\]
assintoticamente quase certamente. Em particular,
\[
  p_H^{\mathrm{rb}}
  \leq
  n^{-1/m_2(H)}
\]
em ordem de grandeza
\cite{KohayakawaKonstadinidisMota2014}.

Assim como no problema Ramsey clássico, a principal dificuldade está no
$0$-enunciado: determinar em quais situações é possível construir uma
coloração própria de $G(n,p)$ sem uma cópia arco-íris de $H$ quando
$p$ está abaixo da escala sugerida pela $2$-densidade.

\subsubsection*{O problema determinístico associado}

Um dos principais avanços nessa direção foi o framework desenvolvido por
Nenadov, Person, Škorić e Steger
\cite{NenadovPersonSkoricSteger2017}. A filosofia desse método consiste
em transformar o problema probabilístico em uma questão determinística
sobre grafos de densidade limitada.

Para um grafo $G$, recordemos que
\[
  m(G)
  :=
  \max_{\substack{J\subseteq G\\v(J)\geq1}}
  \frac{e(J)}{v(J)}.
\]
Quando $H$ é estritamente $2$-balanceado, uma condição natural é
\[
  m(G)\leq m_2(H)
  \quad\Longrightarrow\quad
  G\not\xrightarrow{\mathrm{rb}}H.
  \tag{$\mathcal D_H$}
  \label{eq:anti-condition}
\]
Em outras palavras, todo grafo cuja densidade máxima é, no máximo, a
$2$-densidade do alvo deve admitir uma coloração própria sem uma cópia
arco-íris de $H$.

Sob as hipóteses do framework, a validade de
\eqref{eq:anti-condition} fornece o $0$-enunciado na escala
\[
  n^{-1/m_2(H)}.
\]
Dessa forma, uma questão probabilística sobre $G(n,p)$ é reduzida a um
problema puramente determinístico de coloração.

Neste projeto estaremos inicialmente interessados no limiar em ordem de
grandeza, isto é, em resultados no regime
\[
  p\ll n^{-1/m_2(H)}.
\]
Nesse caso, a condição determinística que aparecerá naturalmente em
muitas aplicações será a versão estrita
\[
  m(G)<m_2(H)
  \quad\Longrightarrow\quad
  G\not\xrightarrow{\mathrm{rb}}H.
  \tag{$\mathcal D_H^{<}$}
  \label{eq:anti-condition-strict}
\]
A distinção entre as versões estrita e não estrita será importante
quando se pretender posteriormente fortalecer os resultados para
limiares semi-agudos.

\subsubsection*{Resultados conhecidos}

O problema apresenta comportamentos bastante distintos dependendo da
estrutura do grafo $H$. O framework de
Nenadov, Person, Škorić e Steger resolveu inicialmente o $0$-enunciado
na escala $n^{-1/m_2(H)}$ para ciclos suficientemente longos e cliques
suficientemente grandes.

Posteriormente, esses resultados foram refinados. Para cliques,
Kohayakawa, Mota, Parczyk e Schnitzer determinaram que
\[
  p_{K_\ell}^{\mathrm{rb}}
  =
  n^{-1/m_2(K_\ell)}
  =
  n^{-2/(\ell+1)}
\]
para todo $\ell\geq5$
\cite{KohayakawaMotaParczykSchnitzer2023}. Entretanto,
\[
  p_{K_4}^{\mathrm{rb}}
  =
  n^{-7/15},
\]
que é assintoticamente menor que
\[
  n^{-1/m_2(K_4)}
  =
  n^{-2/5}.
\]

Um fenômeno análogo ocorre para ciclos. Barros, Cavalar, Mota e Parczyk
provaram que
\[
  p_{C_\ell}^{\mathrm{rb}}
  =
  n^{-1/m_2(C_\ell)}
  =
  n^{-(\ell-2)/(\ell-1)}
\]
para todo $\ell\geq5$, enquanto
\[
  p_{C_4}^{\mathrm{rb}}
  =
  n^{-3/4}
  \ll
  n^{-1/m_2(C_4)}
  =
  n^{-2/3}
\]
\cite{BarrosCavalarMotaParczyk2022}.

Esses exemplos mostram que a escala $n^{-1/m_2(H)}$ não governa
universalmente o problema anti-Ramsey. De fato, são conhecidas
infinitas famílias de grafos para as quais o limiar ocorre
assintoticamente abaixo dessa escala. Entre as obstruções conhecidas
aparecem grafos contendo triângulos, além de exemplos pequenos como
$K_4$, $C_4$ e $K_4$ menos uma aresta.

Por outro lado, resultados recentes mostram que o comportamento previsto
pela $2$-densidade volta a ser universal quando $H$ é suficientemente
denso. Kuperwasser provou que
\[
  m_2(H)\geq19
\]
é suficiente para que
\[
  p_H^{\mathrm{rb}}
  =
  n^{-1/m_2(H)}
\]
em ordem de grandeza \cite{Kuperwasser2026}. A prova é baseada em uma
decomposição de grafos inspirada no teorema das Nine Dragon Trees e
fornece, em particular, uma coloração própria na qual todo subgrafo
arco-íris possui $2$-densidade estritamente menor que um parâmetro de
densidade do grafo ambiente.

Há ainda resultados positivos para outras classes de grafos
estritamente $2$-balanceados, obtidos no contexto de Ramsey restrito
por Behague, Hancock, Hyde, Letzter e Morrison
\cite{BehagueHancockHydeLetzterMorrison2025}. Esses resultados mostram
que a validade da condição determinística não depende exclusivamente
de uma hipótese de alta $2$-densidade e reforçam a importância de se
buscar propriedades estruturais do grafo-alvo.

\subsubsection*{Grafos de cintura grande}

Os exemplos conhecidos sugerem investigar uma direção complementar à de
alta densidade. Enquanto muitas das obstruções conhecidas apresentam
triângulos ou ciclos de comprimento quatro, grafos de cintura grande
possuem uma estrutura local muito mais restrita.

Isso motiva a seguinte questão, que será um dos problemas centrais deste
projeto.

\begin{question}
  \label{que:antiramsey-girth}
  Seja $H$ um grafo que contém um ciclo e satisfaz
  \[
    \operatorname{girth}(H)>4.
  \]
  É verdade que
  \[
    p_H^{\mathrm{rb}}
    =
    n^{-1/m_2(H)}
  \]
  em ordem de grandeza?
\end{question}

Em vista da redução descrita anteriormente, para grafos estritamente
$2$-balanceados uma primeira versão determinística da
Questão~\ref{que:antiramsey-girth} é a seguinte.

\begin{problem}
\label{prob:antiramsey-girth-deterministic}
Determine condições de cintura sobre $H$ que garantam que todo grafo
$G$ satisfazendo
\[
  m(G)<m_2(H)
\]
admite uma coloração própria de suas arestas sem uma cópia
arco-íris de $H$.
\end{problem}

Uma estratégia será começar com uma hipótese de cintura mais forte e,
a partir da compreensão das estruturas que aparecem na prova, tentar
reduzi-la progressivamente. Em particular, uma primeira meta será
investigar o Problema~\ref{prob:antiramsey-girth-deterministic} sob a
hipótese
\[
  \operatorname{girth}(H)>8,
\]
e posteriormente estudar se as mesmas ideias podem ser refinadas até a
condição conjectural
\[
  \operatorname{girth}(H)>4.
\]

Essa abordagem possui a vantagem de separar claramente as duas partes
do problema. A componente probabilística é tratada por ferramentas já
desenvolvidas para $0$-enunciados Ramsey, enquanto a principal questão
nova passa a ser encontrar uma coloração determinística adequada para
grafos de densidade limitada.

\subsubsection*{Estratégias de coloração}

Uma primeira abordagem que pretendemos explorar combina decomposições
por degenerescência com a hipótese de cintura. A ideia é decompor um
grafo $G$ com
\[
  m(G)<m_2(H)
\]
em partes estruturalmente simples e colorir essas partes de forma
coordenada.

Um modelo importante para essa estratégia aparece no trabalho recente
de Kuperwasser. Em linhas gerais, procura-se um subgrafo
$B\subseteq G$ de grau máximo limitado com a propriedade de que toda
cópia de $H$ utilize muitas arestas de $B$. As arestas de $B$ podem
então ser coloridas propriamente com poucas cores, enquanto as arestas
de $G\setminus B$ recebem cores novas. Se toda cópia de $H$ utiliza
mais arestas de $B$ do que o número de cores empregado em $B$, então
alguma cor necessariamente se repete nessa cópia, impedindo que ela
seja arco-íris.

No contexto deste projeto, pretendemos investigar se a ausência de
ciclos curtos em $H$ permite obter versões mais eficientes dessa
decomposição. A hipótese de cintura pode restringir a maneira como uma
cópia de $H$ se distribui entre uma parte degenerada e uma parte de
grau limitado, permitindo reduzir significativamente os parâmetros
necessários no argumento geral.

Uma segunda possibilidade é utilizar decomposições em florestas.
Parâmetros como arboricidade e degenerescência fornecem partições
naturais das arestas de grafos de densidade limitada. Pretendemos
investigar colorações próprias dessas florestas nas quais cores possam
ser reutilizadas de maneira controlada, de forma que toda cópia de
$H$ contenha necessariamente duas arestas de mesma cor. Novamente, a
cintura de $H$ deverá desempenhar um papel importante no controle das
possíveis interseções entre as diferentes partes da decomposição.

Também será considerada uma abordagem indutiva por contraexemplo
minimal. Nesse cenário, assume-se a existência de um grafo $G$ de
densidade mínima possível que seja anti-Ramsey para $H$ e procura-se
remover um vértice de baixo grau, colorir o grafo restante e estender
a coloração. A condição
\[
  m(G)<m_2(H)
\]
fornece controle sobre a degenerescência de $G$, enquanto a hipótese
de cintura restringe as configurações locais que poderiam impedir a
extensão da coloração.

\subsubsection*{Problemas intermediários e extensões}

Além da Questão~\ref{que:antiramsey-girth}, pretendemos investigar
alguns problemas intermediários que podem fornecer resultados
independentes e, ao mesmo tempo, esclarecer a estrutura da questão
principal.

\begin{enumerate}[label=\textup{(AR\arabic*)}]
  \item
        Determinar uma constante $g_0$ para a qual a condição
        \[
          \operatorname{girth}(H)>g_0
        \]
        implique a validade de
        \eqref{eq:anti-condition-strict} para todo grafo
        estritamente $2$-balanceado $H$.

  \item
        Reduzir progressivamente a hipótese sobre a cintura, tendo como
        objetivo final o caso
        \[
          \operatorname{girth}(H)>4.
        \]

  \item
        Investigar separadamente classes de grafos de cintura grande que
        apresentem estrutura adicional, como grafos regulares, quase
        regulares ou bipartidos, comparando os resultados obtidos com as
        classes já cobertas pela literatura.

  \item
        Estender os resultados obtidos para grafos não necessariamente
        $2$-balanceados, analisando o papel dos subgrafos que realizam
        $m_2(H)$.

  \item
        Estudar as obstruções determinísticas de baixa densidade para a
        propriedade anti-Ramsey, buscando identificar quais estruturas
        locais permitem que um grafo $G$ com
        \[
          m(G)<m_2(H)
        \]
        seja anti-Ramsey para $H$.

  \item
        Comparar as técnicas desenvolvidas com variantes do problema em
        que as colorações são localmente limitadas, procurando determinar
        quais argumentos permanecem válidos no contexto Ramsey restrito.
\end{enumerate}

Os problemas acima oferecem diferentes níveis de dificuldade. Os primeiros
resultados poderão ser obtidos sob hipóteses estruturais mais fortes sobre
$H$, enquanto a Questão~\ref{que:antiramsey-girth} representa o objetivo
mais amplo dessa linha de pesquisa. Mesmo que a condição de cintura maior
que $4$ não seja suficiente em toda generalidade, a identificação de novas
famílias positivas ou de novas obstruções de baixa densidade já contribuirá
para uma compreensão mais estrutural do limiar anti-Ramsey.

\subsection{Ramsey canônico assimétrico}
\label{sec:ramsey-canonico-assimetrico}

A segunda linha de pesquisa deste projeto consiste no estudo de uma
versão assimétrica da propriedade Ramsey canônica em grafos aleatórios.
A motivação vem de duas direções clássicas da Teoria de Ramsey:
problemas off-diagonal, nos quais diferentes configurações possuem
diferentes grafos-alvo, e o teorema canônico de Erdős--Rado, no qual
se permitem colorações das arestas com um número arbitrário de cores.

Recordemos que, no problema canônico usual, dada uma ordem total sobre
os vértices, procura-se uma cópia de um grafo $H$ cuja coloração seja
de um dos tipos canônicos: monocromática, arco-íris, min-lex ou
max-lex. Para cliques, Kamčev e Schacht provaram recentemente que,
para todo $\ell\geq4$, o limiar dessa propriedade em $G(n,p)$
coincide com o limiar Ramsey clássico e está na escala
\[
  n^{-1/m_2(K_\ell)}
  =
  n^{-2/(\ell+1)}
\]
\cite{KamcevSchacht2025}.

No contexto determinístico, variantes assimétricas do teorema
canônico também vêm sendo estudadas. Em particular, Araujo e Peng
consideram números off-diagonal de Erdős--Rado nos quais se procuram
cliques de ordens distintas nas configurações monocromática,
lexicográfica e arco-íris
\cite{AraujoPeng2025}. Esses resultados indicam que a assimetria entre
os três padrões canônicos constitui uma extensão natural da teoria
clássica.

Neste projeto propomos estudar a versão probabilística dessa questão.

\begin{definition}
  \label{def:asymmetric-canonical}
  Sejam $H_1,H_2,H_3$ grafos e seja $G$ um grafo equipado com uma
  ordem total $\sigma$ sobre seus vértices. Escrevemos
  \[
    G
    \xrightarrow{\mathrm{can}}
    (H_1,H_2,H_3)
  \]
  quando toda coloração
  \[
    \chi:E(G)\longrightarrow\mathbb N
  \]
  contém ao menos uma das seguintes configurações:
  \begin{enumerate}[label=\textup{(\roman*)}]
    \item uma cópia monocromática de $H_1$;
    \item uma cópia arco-íris de $H_2$;
    \item uma cópia lexicográfica de $H_3$.
  \end{enumerate}
\end{definition}

Nesta formulação inicial, por cópia lexicográfica entendemos uma
cópia min-lex ou max-lex com respeito à ordem induzida por $\sigma$.
Versões em que apenas um desses dois padrões é permitido também são
naturais e poderão ser estudadas separadamente.

Denotaremos por
\[
  \widehat p_{H_1,H_2,H_3}(n)
\]
o limiar correspondente para a propriedade
\[
  G(n,p)
  \xrightarrow{\mathrm{can}}
  (H_1,H_2,H_3).
\]

O primeiro objetivo será compreender essa função quando os três
grafos são completos, isto é,
\[
  (H_1,H_2,H_3)
  =
  (K_r,K_s,K_t).
\]
Mesmo essa família já apresenta diferentes regimes de comportamento,
dependendo de qual dos três parâmetros $r,s,t$ é dominante.

\subsubsection*{Uma primeira cota superior}

Seja
\[
  M:=\max\{r,s,t\}.
\]
Para $M\geq4$, o teorema de Kamčev--Schacht fornece imediatamente uma
cota superior natural para o problema assimétrico. De fato, se
\[
  p\gg n^{-1/m_2(K_M)},
\]
então, assintoticamente quase certamente, toda coloração de $G(n,p)$
contém uma cópia canônica de $K_M$.

Se essa cópia é monocromática, ela contém um $K_r$ monocromático.
Se é arco-íris, contém um $K_s$ arco-íris. Finalmente, se é
lexicográfica, contém um $K_t$ com o mesmo padrão lexicográfico.
Assim,
\[
  \widehat p_{K_r,K_s,K_t}(n)
  \leq
  n^{-1/m_2(K_M)}
  =
  n^{-2/(M+1)}
\]
em ordem de grandeza.

Essa observação sugere investigar quando o maior dos três cliques
também determina o $0$-enunciado.

\begin{question}
  \label{que:canonical-cliques-max}
  Para quais triplas de inteiros $r,s,t\geq3$ vale
  \[
    \widehat p_{K_r,K_s,K_t}(n)
    =
    n^{-1/m_2(K_M)},
    \qquad
    M=\max\{r,s,t\}?
  \]
\end{question}

Resultados preliminares e as diferentes teorias já existentes
sugerem que a resposta depende não apenas do valor de $M$, mas
também da posição ocupada pelo maior clique na tripla.

\subsubsection*{Regimes governados por resultados conhecidos}

Alguns regimes podem ser relacionados diretamente com teoremas
Ramsey ou anti-Ramsey já conhecidos.

Suponha, por exemplo, que o maior clique apareça na posição
monocromática. Uma $2$-coloração sem $K_r$ monocromático evita
automaticamente qualquer clique arco-íris e, quando o alvo
lexicográfico possui ordem pelo menos $4$, também não pode conter uma
cópia lexicográfica estrita, que utiliza pelo menos três cores.
Assim, em diversos casos a cota inferior para o problema assimétrico
pode ser herdada diretamente do $0$-enunciado de Rödl--Ruciński.

De maneira análoga, suponha que o maior clique apareça na posição
arco-íris. Para cliques $K_s$ com $s\geq5$, o limiar anti-Ramsey
próprio é
\[
  n^{-1/m_2(K_s)}.
\]
Abaixo dessa escala existe uma coloração própria sem uma cópia
arco-íris de $K_s$. Como uma coloração própria não possui cliques
monocromáticos de ordem pelo menos $3$ e também não possui cópias
lexicográficas de ordem pelo menos $3$, a mesma coloração pode ser
utilizada diretamente no problema canônico assimétrico.

Essas observações mostram que uma parte da classificação das triplas
de cliques deverá seguir de resultados já conhecidos. O principal
interesse estará nos regimes em que nenhuma dessas reduções imediatas
é suficiente.

\subsubsection*{O caso em que o alvo lexicográfico é dominante}

Uma situação particularmente interessante ocorre quando
\[
  t=\max\{r,s,t\}.
\]
Nesse caso, a escala sugerida é
\[
  n^{-1/m_2(K_t)}
  =
  n^{-2/(t+1)},
\]
mas nem o teorema Ramsey usual nem o resultado anti-Ramsey fornecem
diretamente uma coloração para o $0$-enunciado.

Resultados preliminares obtidos durante o desenvolvimento deste
projeto sugerem uma estratégia baseada nas vizinhanças posteriores
dos vértices.

Fixada uma ordem total $\sigma$ sobre $V(G)$, definimos, para cada
$v\in V(G)$,
\[
  N_\sigma^+(v)
  :=
  \{u\in N_G(v):v<_\sigma u\}
\]
e
\[
  L_v:=G[N_\sigma^+(v)].
\]
A ideia é construir uma coloração das arestas na forma
\[
  \chi(vu)
  =
  \bigl(\lambda(v),\varphi_v(u)\bigr),
  \qquad
  v<_\sigma u,
\]
onde $\lambda$ é injetiva e
\[
  \varphi_v:N_\sigma^+(v)\longrightarrow[q]
\]
é uma coloração auxiliar dos vértices de $L_v$.

A primeira coordenada impede imediatamente cópias monocromáticas de
qualquer clique de ordem pelo menos $3$ e também elimina padrões
max-lex. A segunda coordenada pode então ser escolhida para cumprir
duas funções simultâneas: utilizar poucas cores nas arestas incidentes
ao menor vértice de um clique, impedindo que ele seja arco-íris, e
evitar grandes conjuntos monocromáticos em $L_v$, destruindo os
padrões min-lex.

Para impedir um $K_s$ arco-íris, uma escolha natural é
\[
  q=s-2.
\]
Com essa escolha, as $s-1$ arestas incidentes ao menor vértice de
qualquer $K_s$ utilizam apenas $s-2$ possíveis valores na segunda
coordenada, e portanto duas delas recebem a mesma cor.

Por outro lado, um $K_t$ min-lex com menor vértice $v$ implicaria a
existência de um $K_{t-1}$ monocromático sob $\varphi_v$ em $L_v$.
O problema é então reduzido à existência de uma coloração
\[
  \varphi_v:
  V(L_v)\longrightarrow[s-2]
\]
sem um $K_{t-1}$ monocromático.

Essa redução conecta o problema canônico assimétrico a resultados
sobre densidade Ramsey por vértices e sugere que comparações entre
\[
  m_2(K_t)
\]
e a densidade necessária para forçar
\[
  L_v\xrightarrow{v}(K_{t-1})_{s-2}
\]
podem fornecer colorações para amplas famílias de parâmetros.

\begin{problem}
\label{prob:rst-lex-largest}
Determine para quais $r,s,t$, com $t\geq r,s$, todo grafo
ordenado $B$ satisfazendo
\[
  m(B)<m_2(K_t)
\]
admite uma coloração sem $K_r$ monocromático, sem $K_s$
arco-íris e sem $K_t$ lexicográfico.
\end{problem}

Uma primeira família de interesse é
\[
  3\leq r<s<t.
\]
Os cálculos preliminares indicam que, nesse regime, a estratégia das
vizinhanças é particularmente eficiente. Além disso, a família
consecutiva
\[
  (K_{\ell-1},K_\ell,K_{\ell+1})
\]
constitui um primeiro modelo natural para o problema.

\subsubsection*{Casos pequenos e fenômenos excepcionais}

Assim como no problema anti-Ramsey, os primeiros cliques apresentam
fenômenos especiais que não são necessariamente capturados pelas
construções gerais. Em particular, $K_3$ e $K_4$ desempenham um papel
importante.

Pretendemos estudar inicialmente, entre outros, os casos
\[
  (K_3,K_4,K_3),\qquad
  (K_3,K_3,K_4),\qquad
  (K_3,K_4,K_4),
\]
\[
  (K_4,K_4,K_5)
  \qquad\text{e}\qquad
  (K_3,K_4,K_5).
\]

Esses exemplos são suficientemente pequenos para permitir construções
explícitas e, ao mesmo tempo, ilustram diferentes mecanismos de
coloração.

Resultados preliminares sugerem, por exemplo, que grafos de densidade
máxima menor que
\[
  m_2(K_4)=\frac52
\]
admitem colorações particularmente simples para os casos
$(K_3,K_4,K_3)$ e $(K_3,K_3,K_4)$.

No primeiro caso, a condição de densidade implica degenerescência
limitada e permite colorir propriamente os vértices com cinco cores.
Identificando essas cores com $\mathbb Z_5$, pode-se definir a cor de
uma aresta $uv$ por
\[
  \chi(uv)
  =
  \kappa(u)+\kappa(v)
  \pmod 5.
\]
Essa construção torna todos os triângulos arco-íris e, por utilizar
apenas cinco cores, impede cópias arco-íris de $K_4$.

No caso $(K_3,K_3,K_4)$, uma coloração própria dos vértices com cinco
cores pode ser combinada com uma $2$-coloração de $K_5$ sem triângulo
monocromático. A coloração induzida nas arestas utiliza apenas duas
cores, não possui triângulos monocromáticos e, consequentemente,
também não possui triângulos arco-íris ou cópias lexicográficas
estritas de $K_4$.

Essas construções ainda deverão ser integradas cuidadosamente à parte
probabilística do argumento, mas constituem evidência de que pequenas
obstruções podem ser tratadas por \emph{templates} finitos de
coloração. Essa técnica deverá ser explorada sistematicamente ao longo
do projeto.

\subsubsection*{Templates finitos e degenerescência}

Mais geralmente, suponha que a condição
\[
  m(B)<d
\]
implique que $B$ seja $q$-colorível nos vértices. Uma estratégia
natural consiste em escolher uma coloração das arestas de um grafo
completo $K_q$ que já evite os padrões desejados e transferi-la para
$B$ através de uma coloração própria
\[
  \kappa:V(B)\longrightarrow V(K_q).
\]

Mais precisamente, se
\[
  \rho:E(K_q)\longrightarrow\mathcal C
\]
é uma coloração adequada, definimos
\[
  \chi(uv)
  :=
  \rho\bigl(\kappa(u)\kappa(v)\bigr).
\]

Como $\kappa$ é injetiva em todo clique de $B$, cada clique de $B$
reproduz exatamente o padrão de algum clique do template $K_q$.
Assim, problemas infinitos de coloração em todos os grafos de
determinada densidade podem, em certos casos, ser reduzidos à procura
de um template finito.

Pretendemos investigar quais triplas $(r,s,t)$ admitem esse tipo de
construção e em que situações a busca por templates pode ser
automatizada computacionalmente para sugerir novas colorações e
obstruções. Os resultados computacionais serão utilizados apenas como
ferramenta exploratória; os resultados finais deverão ser acompanhados
de demonstrações matemáticas.

\subsubsection*{Um possível parâmetro de densidade assimétrico}

No problema Ramsey assimétrico clássico, quando se buscam diferentes
grafos em diferentes cores, o limiar não é necessariamente governado
simplesmente por
\[
  \max_i m_2(H_i).
\]
A conjectura de Kohayakawa--Kreuter introduz densidades mistas que
capturam a interação entre grafos-alvo distintos.

Um dos objetivos de longo prazo desta linha será investigar se existe
um fenômeno semelhante na versão canônica assimétrica.

\begin{question}
  \label{que:mixed-canonical-density}
  Existe um parâmetro
  \[
    m_{\mathrm{can}}(H_1,H_2,H_3)
  \]
  definido em termos das estruturas de $H_1,H_2,H_3$ que determine,
  em uma classe ampla de casos, o limiar
  \[
    \widehat p_{H_1,H_2,H_3}(n)
    =
    n^{-1/m_{\mathrm{can}}(H_1,H_2,H_3)}?
  \]
\end{question}

Para cliques, uma primeira possibilidade a ser testada é que, em
regimes suficientemente regulares,
\[
  m_{\mathrm{can}}(K_r,K_s,K_t)
  =
  m_2(K_{\max\{r,s,t\}}).
\]
Os casos pequenos deverão indicar precisamente onde essa heurística
falha e quais informações adicionais precisam entrar no parâmetro.

Outra possibilidade é que diferentes posições da tripla sejam
governadas por parâmetros de naturezas distintas: densidades Ramsey
para a alternativa monocromática, densidades anti-Ramsey para a
alternativa arco-íris e densidades associadas a problemas Ramsey por
vértices para a alternativa lexicográfica. Compreender essa interação
é uma das questões conceituais mais interessantes da proposta.

\subsubsection*{Uma redução probabilístico-determinística}

Um objetivo central será desenvolver uma versão do paradigma de
Nenadov--Person--Škorić--Steger adaptada ao problema canônico
assimétrico.

Idealmente, gostaríamos de obter um princípio da seguinte forma.
Fixados $H_1,H_2,H_3$ e um parâmetro crítico $d$, se todo grafo
ordenado $B$ com
\[
  m(B)<d
\]
admite uma coloração evitando simultaneamente $H_1$ monocromático,
$H_2$ arco-íris e $H_3$ lexicográfico, então
\[
  p\ll n^{-1/d}
\]
deveria permitir construir, assintoticamente quase certamente, uma
coloração correspondente de $G(n,p)$.

A principal dificuldade está em determinar qual das três
configurações deve governar o procedimento de remoção e a definição
dos blocos. Quando o maior alvo é, por exemplo, $K_t$ na posição
lexicográfica, é natural explorar blocos construídos a partir de
cópias de $K_t$. Em outros regimes, o alvo relevante pode estar na
posição monocromática ou arco-íris, e uma decomposição diferente pode
ser mais apropriada.

\begin{problem}
\label{prob:canonical-framework}
Desenvolver condições suficientes sob as quais um $0$-enunciado
para a propriedade
\[
  G(n,p)
  \xrightarrow{\mathrm{can}}
  (H_1,H_2,H_3)
\]
possa ser reduzido a um problema determinístico de coloração em
grafos de tamanho e densidade limitados.
\end{problem}

A obtenção de um resultado desse tipo permitiria separar de maneira
clara a parte probabilística da parte de coloração e tornaria possível
estudar diferentes triplas utilizando um mesmo framework.

\subsubsection*{Problemas propostos}

A investigação será inicialmente organizada em torno dos seguintes
problemas.

\begin{enumerate}[label=\textup{(RC\arabic*)}]
  \item
        Determinar o limiar para as triplas
        \[
          (K_r,K_s,K_t)
        \]
        em todos os regimes possíveis de $r,s,t$, começando pelos casos
        pequenos e pelas famílias em que um dos parâmetros domina os
        demais.

  \item
        Estudar sistematicamente o regime
        \[
          3\leq r<s<t,
        \]
        com particular atenção à família consecutiva
        \[
          (K_{\ell-1},K_\ell,K_{\ell+1}).
        \]

  \item
        Classificar os casos em que o limiar segue diretamente de um
        resultado Ramsey clássico ou anti-Ramsey e identificar os regimes
        genuinamente canônicos.

  \item
        Resolver os casos pequenos envolvendo $K_3$ e $K_4$, incluindo
        \[
          (K_3,K_4,K_3),\quad
          (K_3,K_3,K_4),\quad
          (K_3,K_4,K_4),\quad
          (K_4,K_4,K_5)
        \]
        e outras triplas próximas.

  \item
        Desenvolver e sistematizar métodos de coloração baseados em
        vizinhanças posteriores, degenerescência e templates finitos.

  \item
        Investigar a existência de um parâmetro de densidade misto para
        o problema Ramsey canônico assimétrico.

  \item
        Obter uma redução probabilístico-determinística que permita
        transferir lemas de coloração para $0$-enunciados em $G(n,p)$.

  \item
        Após a compreensão do caso de cliques, estudar triplas envolvendo
        ciclos e outras classes de grafos, como
        \[
          (K_r,C_s,K_t),\qquad
          (C_r,K_s,K_t)
          \qquad\text{e}\qquad
          (C_r,C_s,C_t).
        \]
\end{enumerate}

O estudo das triplas de cliques fornecerá um laboratório para a
formulação de conjecturas gerais. Além de permitir resultados
concretos em estágios iniciais do doutorado, esses casos deverão
indicar quais parâmetros de densidade e quais mecanismos de coloração
são realmente relevantes para a teoria assimétrica.

Em uma etapa posterior, pretendemos utilizar essa compreensão para
investigar grafos menos homogêneos, nos quais diferentes subgrafos
podem realizar as densidades críticas e nos quais a ordem dos vértices
pode influenciar de maneira mais significativa a existência de padrões
lexicográficos.

\subsection{Ramsey canônico para grafos gerais}
\label{sec:ramsey-canonico-geral}

A terceira linha de pesquisa deste projeto busca compreender a propriedade
Ramsey canônica em grafos aleatórios quando o grafo-alvo não é
necessariamente completo.

Seja $H$ um grafo equipado com uma ordem total $\sigma$ sobre seus
vértices. Dado um grafo ordenado $G$, escrevemos
\[
  G\xrightarrow{\mathrm{can}}H
\]
quando toda coloração
\[
  \chi:E(G)\longrightarrow\mathbb N
\]
contém uma cópia de $H$ cuja coloração é canônica, isto é, uma cópia
monocromática, arco-íris, min-lex ou max-lex, respeitando a ordem
prescrita em $H$.

Denotaremos por
\[
  \widehat p_H(n)
\]
o limiar para essa propriedade em $G(n,p)$, equipado com uma ordem
total fixa sobre seus vértices.

No caso em que $H$ é um clique, esse problema foi recentemente
resolvido por Kamčev e Schacht. Para todo $\ell\geq4$,
\[
  \widehat p_{K_\ell}(n)
  =
  n^{-1/m_2(K_\ell)}
  =
  n^{-2/(\ell+1)}
\]
em ordem de grandeza
\cite{KamcevSchacht2025}.

É notável que essa é exatamente a mesma escala que aparece no Teorema
de Ramsey aleatório de Rödl e Ruciński. Assim, ao menos para cliques,
permitir um número arbitrário de cores e substituir a conclusão
monocromática por uma conclusão canônica não altera a ordem do limiar.

Esse fenômeno motiva uma das questões gerais desta linha de pesquisa.

\begin{question}
  \label{que:canonical-general}
  Para quais grafos ordenados $H$ vale
  \[
    \widehat p_H(n)
    =
    n^{-1/m_2(H)}
  \]
  em ordem de grandeza?
\end{question}

A Questão~\ref{que:canonical-general} contém duas dificuldades
distintas. A primeira é probabilística: compreender em que escala
$G(n,p)$ possui estrutura suficiente para forçar um padrão canônico.
A segunda é estrutural e está relacionada à ordenação dos vértices.
Para grafos que não são completos, diferentes ordenações de um mesmo
grafo não precisam se comportar da mesma forma do ponto de vista
canônico.

\subsubsection*{Resultados além dos cliques}

O trabalho de Kamčev e Schacht fornece algumas indicações de que a
escala
\[
  n^{-1/m_2(H)}
\]
deve permanecer relevante para grafos mais gerais. Em particular,
para um grafo estritamente $2$-balanceado $H$, suas técnicas garantem,
no regime
\[
  p\gg n^{-1/m_2(H)},
\]
a existência de uma cópia canonicamente colorida de $H$. Entretanto,
nessa formulação não se pode, em geral, prescrever previamente a
ordenação dos vértices da cópia.

Isso revela uma distinção que é invisível no caso de cliques. Como
toda ordenação de $K_\ell$ é isomorfa como grafo ordenado, a ordem não
representa uma dificuldade adicional. Para um grafo geral $H$, porém,
o problema de encontrar alguma ordenação canônica e o problema de
encontrar uma ordenação previamente fixada podem ser genuinamente
diferentes.

Outra classe para a qual já existem avanços importantes é a dos ciclos
pares. Alvarado, Kohayakawa, Morris e Mota provaram um teorema
canônico para $C_{2\ell}$ em grafos aleatórios
\cite{AlvaradoKohayakawaMorrisMotaCycles}. Em particular, seu resultado
fornece uma cópia canônica quando
\[
  p
  \gg
  n^{-1/m_2(C_{2\ell})}\log n.
\]
Como
\[
  m_2(C_{2\ell})
  =
  \frac{2\ell-1}{2\ell-2},
\]
essa escala pode ser escrita como
\[
  p
  \gg
  n^{-1+1/(2\ell-1)}\log n.
\]

O resultado determina, portanto, o limiar esperado até um fator
logarítmico. Há, contudo, uma diferença importante entre essa
formulação e a versão ordenada que pretendemos estudar neste projeto:
no resultado para ciclos pares, a ordenação relevante pode ser
escolhida após a coloração do grafo aleatório.

Esses resultados sugerem que há pelo menos três versões naturais do
problema:
\begin{enumerate}[label=\textup{(\roman*)}]
  \item encontrar alguma ordenação de $H$ cuja cópia seja canônica;
  \item fixar a ordem do grafo ambiente antes da coloração, mas
        permitir qualquer tipo ordenado de $H$;
  \item fixar previamente um grafo ordenado $(H,\sigma)$ e exigir
        uma cópia que preserve exatamente essa ordem.
\end{enumerate}

Pretendemos compreender em que situações essas versões possuem o
mesmo limiar e em que situações a ordem prescrita introduz um novo
fenômeno.

\subsubsection*{Ciclos}

Os ciclos constituem uma primeira classe natural para investigar a
Questão~\ref{que:canonical-general}. Além de possuírem uma estrutura
simples e $2$-densidade explicitamente conhecida, eles já apresentam
uma diferença fundamental entre os casos par e ímpar.

Para todo $\ell\geq3$,
\[
  m_2(C_\ell)
  =
  \frac{\ell-1}{\ell-2},
\]
de modo que a escala sugerida pelo problema Ramsey clássico é
\[
  n^{-1/m_2(C_\ell)}
  =
  n^{-(\ell-2)/(\ell-1)}.
\]

Para ciclos pares, o primeiro problema será entender se o fator
logarítmico presente nos resultados conhecidos é necessário.

\begin{problem}
\label{prob:even-cycle-canonical}
Determine o limiar exato para a propriedade Ramsey canônica de
$C_{2\ell}$. Em particular, determine se
\[
  \widehat p_{C_{2\ell}}(n)
  =
  n^{-1/m_2(C_{2\ell})}
\]
em ordem de grandeza.
\end{problem}

Uma etapa adicional será estudar a versão com ordenação previamente
prescrita.

\begin{problem}
\label{prob:ordered-even-cycle}
Seja $(C_{2\ell},\sigma)$ um ciclo par equipado com uma ordem total
fixa sobre seus vértices. Determine o limiar para que toda
coloração arbitrária de $G(n,p)$ contenha uma cópia canônica de
$(C_{2\ell},\sigma)$.
\end{problem}

Para ciclos ímpares, a situação parece ainda menos compreendida, o que
motiva uma segunda família de problemas.

\begin{problem}
\label{prob:odd-cycle-canonical}
Para $\ell\geq2$, determine o limiar Ramsey canônico de
\[
  C_{2\ell+1}.
\]
Em particular, determine se a escala correta é
\[
  n^{-1/m_2(C_{2\ell+1})}
  =
  n^{-(2\ell-1)/(2\ell)}.
\]
\end{problem}

O primeiro caso não trivial,
\[
  C_5,
\]
deverá servir como modelo inicial. A análise de $C_5$ pode revelar
fenômenos que não aparecem para cliques e, ao mesmo tempo, permanece
suficientemente pequena para permitir o estudo explícito das possíveis
ordenações e padrões de coloração.

Posteriormente, pretende-se investigar $C_7$ e $C_9$ e buscar uma
construção que se estenda a ciclos ímpares arbitrários.

\subsubsection*{O papel da ordenação}

Um aspecto que pretendemos estudar sistematicamente é a dependência do
limiar em relação à ordenação escolhida para o grafo-alvo.

Para um grafo $H$ e duas ordens totais $\sigma_1,\sigma_2$ sobre
$V(H)$, não é evidente que
\[
  \widehat p_{(H,\sigma_1)}
\]
e
\[
  \widehat p_{(H,\sigma_2)}
\]
devam possuir a mesma ordem de grandeza.

Isso sugere definir os parâmetros
\[
  \widehat p_H^{\min}(n)
  :=
  \min_{\sigma}
  \widehat p_{(H,\sigma)}(n)
\]
e
\[
  \widehat p_H^{\max}(n)
  :=
  \max_{\sigma}
  \widehat p_{(H,\sigma)}(n)
\]
e investigar quão diferentes essas funções podem ser.

\begin{question}
  \label{que:canonical-order}
  Existe um grafo $H$ para o qual duas ordenações $\sigma_1$ e
  $\sigma_2$ produzem limiares Ramsey canônicos de ordens de
  grandeza distintas?
\end{question}

Mesmo que os limiares sejam sempre da mesma ordem de grandeza, as
constantes ou os mecanismos envolvidos nos $0$-enunciados podem
depender significativamente da ordem.

Para ciclos, essa questão já pode ser estudada concretamente. Uma
ordenação em que os vértices aparecem na ordem cíclica é
estruturalmente bastante diferente de uma ordenação alternada entre
as duas partes de um ciclo par, por exemplo. Pretendemos classificar,
ao menos para ciclos pequenos, os diferentes tipos ordenados e
comparar seus respectivos comportamentos.

\subsubsection*{Grafos estritamente $2$-balanceados}

A evidência existente sugere que a classe dos grafos estritamente
$2$-balanceados é um ponto natural de partida para uma teoria mais
geral.

\begin{question}
  \label{que:strictly-balanced-canonical}
  Seja $H$ um grafo estritamente $2$-balanceado que contém um ciclo.
  É verdade que existe uma ordenação $\sigma$ para a qual
  \[
    \widehat p_{(H,\sigma)}(n)
    =
    n^{-1/m_2(H)}?
  \]
\end{question}

Uma versão mais forte seria exigir o resultado para toda ordenação.

\begin{question}
  \label{que:every-order-canonical}
  Seja $H$ estritamente $2$-balanceado. Para toda ordem total
  $\sigma$ sobre $V(H)$, vale
  \[
    \widehat p_{(H,\sigma)}(n)
    =
    n^{-1/m_2(H)}?
  \]
\end{question}

Os resultados existentes fornecem evidência para a primeira questão,
mas a segunda envolve diretamente a dificuldade adicional de
prescrever a ordem.

Entre as classes a serem consideradas nessa direção estão:
\begin{itemize}
  \item ciclos;
  \item grafos regulares esparsos;
  \item grafos de cintura grande;
  \item grafos bipartidos estritamente $2$-balanceados;
  \item grafos com degenerescência limitada.
\end{itemize}

A escolha dessas classes também permitirá comparar os fenômenos
canônicos com aqueles observados na parte anti-Ramsey deste projeto.

\subsubsection*{$0$-enunciados e problemas determinísticos}

Assim como nas duas linhas anteriores, um objetivo importante será
compreender os $0$-enunciados através de problemas determinísticos de
coloração.

Para cliques, a escala crítica sugere estudar grafos $B$ satisfazendo
\[
  m(B)<m_2(K_\ell)
\]
e construir uma coloração de suas arestas sem uma cópia canônica de
$K_\ell$.

Para grafos gerais, entretanto, o problema determinístico pode
depender fortemente da ordem e da estrutura de $H$. Isso motiva a
seguinte questão.

\begin{problem}
\label{prob:canonical-deterministic-general}
Dado um grafo ordenado $(H,\sigma)$, determine condições sobre
um grafo ordenado $B$ com
\[
  m(B)<m_2(H)
\]
que garantam a existência de uma coloração de $E(B)$ sem uma
cópia canônica de $(H,\sigma)$.
\end{problem}

A ausência de uma cópia canônica significa evitar simultaneamente
cópias monocromáticas, arco-íris, min-lex e max-lex. Diferentes partes
dessa condição sugerem técnicas distintas.

Uma coloração com poucas cores elimina automaticamente cópias
arco-íris grandes. Uma coloração própria elimina configurações
monocromáticas e lexicográficas. Colorações baseadas em
degenerescência ou em templates finitos podem permitir um compromisso
entre esses dois extremos.

Pretendemos investigar em quais classes de grafos essas estratégias
podem ser combinadas e se uma decomposição do grafo ambiente em
partes com diferentes tipos de coloração fornece uma abordagem
sistemática para o problema.

\subsubsection*{Relação com resultados de Ramsey canônico determinístico}

Outra direção será utilizar construções determinísticas Ramsey
canônicas como inspiração para o problema aleatório.

Resultados clássicos de Erdős--Rado e Nešetřil--Rödl mostram que é
possível construir grafos com propriedades Ramsey canônicas bastante
fortes. Trabalhos recentes também investigam grafos Ramsey canônicos
com restrições estruturais adicionais.

Uma questão natural é determinar até que ponto tais construções podem
ser combinadas com argumentos probabilísticos ou de esparsificação.

\begin{question}
  \label{que:sparse-canonical-host}
  Dado um grafo ordenado $H$, quão esparso pode ser um grafo $G$
  satisfazendo
  \[
    G\xrightarrow{\mathrm{can}}H?
  \]
  Em particular, quais parâmetros de densidade são necessários em
  grafos Ramsey canônicos para $H$?
\end{question}

Essa questão determinística está intimamente relacionada ao estudo dos
limiares. A existência de uma obstrução Ramsey canônica $G$ com
\[
  m(G)<m_2(H)
\]
poderia indicar que a escala natural prevista por $m_2(H)$ não é
correta, de maneira semelhante ao que ocorre no problema anti-Ramsey.

Assim, além de provar resultados positivos, será importante procurar
sistematicamente pequenas obstruções que possam produzir exceções.

\subsubsection*{Métodos com listas}

Uma ferramenta adicional que poderá ser utilizada é a versão Ramsey
canônica com restrições por listas desenvolvida por Alvarado,
Kohayakawa, Morris e Mota
\cite{AlvaradoKohayakawaMorrisMotaLists}.

Problemas com listas surgem naturalmente quando uma coloração global é
construída em etapas. Após colorir uma parte do grafo, as cores
disponíveis para as arestas restantes podem ser descritas por listas,
e um resultado canônico para colorações sujeitas a essas listas pode
permitir concluir o argumento.

Essa técnica já desempenha papel importante no estudo de ciclos pares
e poderá ser útil tanto na remoção do fator logarítmico quanto no
tratamento de outras classes de grafos.

\subsubsection*{Problemas propostos}

A terceira linha do projeto será inicialmente organizada em torno dos
seguintes problemas.

\begin{enumerate}[label=\textup{(CG\arabic*)}]
  \item
        Determinar o limiar Ramsey canônico para ciclos ímpares,
        começando por
        \[
          C_5
          \qquad\text{e}\qquad
          C_7.
        \]

  \item
        Determinar se o fator logarítmico presente no resultado conhecido
        para ciclos pares pode ser removido.

  \item
        Estudar o limiar de ciclos pares quando uma ordenação específica
        é fixada previamente.

  \item
        Classificar as diferentes ordenações de ciclos pequenos do ponto
        de vista da propriedade Ramsey canônica.

  \item
        Determinar se diferentes ordenações de um mesmo grafo podem
        produzir limiares de ordens de grandeza distintas.

  \item
        Investigar se todo grafo estritamente $2$-balanceado possui
        limiar canônico governado por
        \[
          n^{-1/m_2(H)}
        \]
        quando não se prescreve a ordem.

  \item
        Investigar a mesma questão quando uma ordenação arbitrária é
        previamente fixada.

  \item
        Desenvolver critérios determinísticos para colorir grafos
        $B$ satisfazendo
        \[
          m(B)<m_2(H)
        \]
        sem produzir cópias canônicas de $H$.

  \item
        Procurar obstruções de baixa densidade que possam produzir
        exceções ao comportamento governado por $m_2(H)$.

  \item
        Estender os resultados obtidos para outras classes naturais,
        incluindo grafos regulares esparsos, grafos bipartidos e grafos
        de cintura grande.

  \item
        Investigar extensões da propriedade canônica para famílias
        finitas de grafos e, em uma etapa posterior, para hipergrafos.
\end{enumerate}

O estudo de ciclos deverá fornecer a primeira ponte entre o resultado
para cliques e uma teoria Ramsey canônica para grafos gerais. Ao mesmo
tempo, a dependência em relação à ordenação introduz um fenômeno que
não aparece no caso completo e que poderá exigir novas ferramentas.

Em uma perspectiva mais ampla, pretende-se compreender se a
$2$-densidade máxima desempenha na Teoria de Ramsey canônica o mesmo
papel universal que desempenha na Teoria de Ramsey aleatória clássica,
ou se existem classes de grafos e ordenações para as quais novos
parâmetros estruturais são necessários.
% ============================================================
% METODOLOGIA DE ESTUDO
% ============================================================

\section{Metodologia de estudo}
\label{sec:metodologia}

Os problemas apresentados na Seção~\ref{sec:anti-ramsey} e nas
subseções seguintes encontram-se na interseção entre Teoria de Ramsey,
Combinatória Extremal e Teoria de Grafos Aleatórios. Seu estudo exigirá
a combinação de técnicas probabilísticas com argumentos determinísticos
de coloração e decomposição de grafos.

Uma característica comum às três linhas propostas é a separação entre
os chamados $0$-enunciados e $1$-enunciados. Para uma propriedade
monótona $\mathcal P$ em $G(n,p)$, o $1$-enunciado consiste em mostrar
que
\[
  \Pr[G(n,p)\in\mathcal P]\longrightarrow1
\]
acima de determinada escala de $p$, enquanto o $0$-enunciado consiste
em construir, abaixo dessa escala, uma coloração que testemunhe a
falha da propriedade.

Embora essas duas partes estejam fortemente relacionadas, as técnicas
necessárias para tratá-las costumam ser diferentes. Resultados de
transferência, métodos de containers e versões esparsas de teoremas
extremais são particularmente apropriados para $1$-enunciados. Por
outro lado, os $0$-enunciados frequentemente exigem uma compreensão
mais fina da estrutura local do grafo aleatório e a construção
explícita de colorações que evitem as configurações desejadas.

Uma das principais estratégias metodológicas deste projeto será
explorar essa distinção, procurando reduzir os $0$-enunciados
probabilísticos a problemas determinísticos de coloração.

\subsection*{Reduções probabilístico-determinísticas}

Um modelo fundamental para essa abordagem é o framework de Nenadov,
Person, Škorić e Steger
\cite{NenadovPersonSkoricSteger2017}. Em diversas propriedades do tipo
Ramsey, o método consiste inicialmente em remover arestas que podem ser
coloridas sem dificuldade e, posteriormente, decompor o subgrafo
remanescente em estruturas menores, usualmente denominadas blocos.

No regime abaixo da escala crítica, esses blocos possuem,
assintoticamente quase certamente, tamanho limitado. Além disso, um
argumento de primeiro momento permite controlar sua densidade. Dessa
forma, a parte probabilística do problema pode ser reduzida a uma
afirmação determinística da forma
\[
  m(B)<d
  \quad\Longrightarrow\quad
  B\text{ admite uma coloração adequada},
\]
onde $d$ é o parâmetro de densidade associado ao problema em questão.

Pretendemos utilizar essa filosofia nas três linhas do projeto.

No problema anti-Ramsey, a questão determinística consiste em produzir
uma coloração própria de $B$ sem uma cópia arco-íris de $H$.

No problema Ramsey canônico assimétrico, será necessário evitar
simultaneamente uma cópia monocromática de $H_1$, uma cópia arco-íris
de $H_2$ e uma cópia lexicográfica de $H_3$.

Finalmente, no problema Ramsey canônico para grafos gerais, será
necessário evitar simultaneamente os diferentes padrões canônicos do
grafo-alvo, levando também em consideração a ordenação prescrita de
seus vértices.

Uma parte importante do projeto será determinar até que ponto o
framework existente pode ser aplicado diretamente a essas propriedades
e quais modificações são necessárias para capturar os diferentes tipos
de coloração envolvidos.

\subsection*{Decomposições de grafos e degenerescência}

Uma segunda família de ferramentas será formada por técnicas de
decomposição de grafos.

Se um grafo $G$ satisfaz uma condição do tipo
\[
  m(G)<d,
\]
então sua densidade limitada impõe restrições à degenerescência, à
arboricidade e à possibilidade de decompor seu conjunto de arestas em
subgrafos estruturalmente simples.

Essas propriedades deverão desempenhar um papel particularmente
importante na linha anti-Ramsey. Uma estratégia será decompor o grafo
ambiente em florestas, grafos de grau máximo controlado ou subgrafos
orientados com grau de saída limitado, procurando reutilizar cores de
forma coordenada. O objetivo será garantir que toda cópia do
grafo-alvo contenha duas arestas de mesma cor, sem violar a condição de
a coloração global ser própria.

A hipótese de cintura grande deverá ser combinada com essas
decomposições. Como duas regiões distintas de um grafo de cintura
grande não podem ser conectadas por caminhos curtos arbitrários,
esperamos que a ausência de ciclos pequenos restrinja a maneira como
uma cópia de $H$ pode se distribuir entre as diferentes partes de uma
decomposição.

Pretendemos também estudar decomposições inspiradas em resultados
recentes utilizados no problema anti-Ramsey, incluindo decomposições
em florestas e subgrafos de grau limitado.

\subsection*{Colorações locais e vizinhanças}

Na linha Ramsey canônica assimétrica, uma técnica que deverá receber
atenção especial é a utilização das vizinhanças orientadas pela ordem
dos vértices.

Dada uma ordem total $\sigma$, para cada vértice $v$ consideramos
\[
  N_\sigma^+(v)
  =
  \{u\in N_G(v):v<_\sigma u\}.
\]
Uma coloração local dos vértices do grafo
\[
  G[N_\sigma^+(v)]
\]
pode ser utilizada para definir as cores das arestas incidentes a
$v$.

Construções desse tipo permitem transformar um problema de coloração
de arestas em um problema Ramsey por vértices. Em particular, podem
ser utilizadas simultaneamente para limitar o número de cores
presentes nas arestas incidentes ao menor vértice de um clique e para
impedir padrões lexicográficos.

Essa abordagem deverá ser comparada com resultados conhecidos sobre
densidade Ramsey por vértices e poderá fornecer construções gerais
para famílias inteiras de triplas
\[
  (K_r,K_s,K_t).
\]

\subsection*{Templates finitos de coloração}

Outra técnica que aparece naturalmente nos casos pequenos consiste no
uso de templates finitos.

Se uma condição de densidade garante que um grafo $B$ admite uma
coloração própria de seus vértices com $q$ cores, podemos escolher uma
coloração auxiliar
\[
  \rho:E(K_q)\longrightarrow\mathcal C
\]
e transferi-la para $B$ através de uma coloração própria
\[
  \kappa:V(B)\longrightarrow[q]
\]
definindo
\[
  \chi(uv)
  :=
  \rho(\kappa(u)\kappa(v)).
\]

Como toda cópia de um clique em $B$ utiliza vértices de cores distintas
sob $\kappa$, seu padrão de coloração é reproduzido por uma cópia do
mesmo clique no template $K_q$.

Essa técnica já fornece construções naturais para alguns dos primeiros
casos da propriedade Ramsey canônica assimétrica e será investigada de
maneira sistemática.

Além de construções explícitas, será possível realizar buscas
computacionais por pequenos templates. Esses experimentos terão caráter
exploratório: poderão sugerir padrões de coloração, identificar casos
excepcionais ou indicar conjecturas, mas os resultados matemáticos do
projeto serão acompanhados de demonstrações independentes das buscas
computacionais.

\subsection*{Método probabilístico e argumentos de contagem}

Ferramentas clássicas do método probabilístico serão utilizadas ao
longo de todo o projeto.

Argumentos de primeiro momento serão empregados, por exemplo, para
mostrar que determinados subgrafos de densidade elevada não aparecem
em $G(n,p)$ abaixo de uma escala crítica. Em situações mais delicadas,
poderão ser utilizados argumentos de segundo momento, desigualdades de
concentração e resultados gerais sobre o aparecimento de subgrafos
fixos em grafos aleatórios.

Para $1$-enunciados, pretendemos utilizar resultados existentes sempre
que possível. Por exemplo, o teorema Ramsey canônico de
Kamčev--Schacht fornece diretamente vários limites superiores para os
problemas estudados neste projeto.

Quando os resultados existentes não puderem ser aplicados
diretamente, estudaremos adaptações de métodos de containers de
hipergrafos e princípios de transferência para ambientes esparsos.

O método de containers é particularmente natural em problemas Ramsey:
uma coloração que evita determinada configuração pode ser descrita
através de conjuntos independentes de um hipergrafo auxiliar cujos
vértices representam as arestas do grafo ambiente e cujas
hiperarestas correspondem às cópias da configuração proibida.

\subsection*{Colorações canônicas com listas}

Na investigação de grafos gerais, especialmente ciclos, pretendemos
também explorar versões da Teoria de Ramsey canônica com listas de
cores.

Essas versões são úteis quando uma coloração é construída
iterativamente. Depois de colorida parte do grafo, as escolhas
possíveis para uma nova aresta podem ser representadas por uma lista de
cores admissíveis.

Resultados Ramsey canônicos com listas poderão, portanto, funcionar
como ferramenta intermediária para construir colorações globais ou
para encontrar padrões canônicos em subgrafos esparsos.

Essa abordagem já aparece em trabalhos recentes relacionados ao
problema canônico para ciclos pares e será estudada com o objetivo de
adaptá-la a ciclos ímpares, ordenações prescritas e outros grafos
esparsos.

\subsection*{Análise de obstruções pequenas}

Além de procurar provas gerais, pretendemos estudar sistematicamente
possíveis obstruções determinísticas.

Dado um problema cuja escala candidata é
\[
  n^{-1/d},
\]
um grafo fixo $B$ satisfazendo
\[
  m(B)<d
\]
e que já possua a propriedade Ramsey correspondente pode impedir que o
limiar conjecturado seja correto.

Esse fenômeno ocorre, por exemplo, em diferentes variantes anti-Ramsey
e mostra a importância de compreender estruturas pequenas antes de
formular conjecturas gerais.

Nos casos envolvendo cliques pequenos e ciclos pequenos, pretende-se
combinar argumentos matemáticos com enumeração computacional de grafos
de ordem limitada. O objetivo será identificar potenciais obstruções,
classificá-las e compreender quais características estruturais são
responsáveis por seu comportamento.

Essa análise poderá também fornecer evidências para a escolha de novos
parâmetros de densidade nos casos em que $m_2(H)$ não seja suficiente.

\subsection*{Desenvolvimento da pesquisa}

O desenvolvimento do projeto será acompanhado por estudo contínuo da
literatura recente em Teoria de Ramsey, grafos aleatórios e
Combinatória Extremal. Parte importante do trabalho inicial será a
consolidação das técnicas presentes nos resultados de
Rödl--Ruciński, Nenadov--Steger, Nenadov--Person--Škorić--Steger,
Kamčev--Schacht e nos trabalhos recentes sobre propriedades
anti-Ramsey.

A pesquisa será desenvolvida em reuniões regulares com o orientador e
o coorientador. Além da discussão dos problemas específicos do
projeto, essas reuniões serão utilizadas para a análise crítica das
estratégias empregadas, escolha das direções mais promissoras e
acompanhamento dos resultados obtidos.

O candidato participará também dos seminários e atividades dos grupos
de Combinatória e Teoria dos Grafos do Instituto de Matemática e
Estatística da Universidade de São Paulo, mantendo contato com
pesquisadores visitantes e colaboradores que trabalham em temas
próximos aos propostos neste projeto.

A interação com outros pesquisadores será particularmente importante,
uma vez que os problemas investigados se relacionam diretamente com
linhas de pesquisa ativas em Teoria de Ramsey aleatória,
anti-Ramsey, métodos probabilísticos e Combinatória Extremal.

Pretende-se ainda realizar, durante o doutorado, um estágio de pesquisa
no exterior em um grupo com atuação em Teoria de Ramsey e grafos
aleatórios. A escolha do grupo e do pesquisador anfitrião será feita
em função do desenvolvimento das linhas propostas e dos problemas que
se mostrarem mais promissores ao longo dos primeiros anos do projeto.

% ============================================================
% O CANDIDATO
% ============================================================

\section{O candidato}
\label{sec:candidato}

O candidato concluiu o mestrado em Ciência da Computação no Instituto de
Matemática e Estatística da Universidade de São Paulo em 2025, sob
orientação de Yoshiharu Kohayakawa. Sua dissertação, intitulada
\emph{Classes Ramsey e Colorações Canônicas via Método Partido}, foi
dedicada ao estudo de problemas em Teoria de Ramsey, com ênfase em classes
Ramsey, hipergrafos, colorações canônicas e no método partido de
Nešetřil--Rödl.

Durante o mestrado, o candidato adquiriu formação em Combinatória Extremal
e Teoria de Ramsey, estudando tanto resultados clássicos quanto técnicas
estruturais utilizadas em problemas modernos da área. Uma parte importante
do trabalho foi dedicada ao método partido, ferramenta introduzida por
Nešetřil e Rödl que permite construir estruturas Ramsey preservando
restrições combinatórias adicionais.

Na primeira parte da dissertação, o candidato estudou em detalhe
propriedades Ramsey de sistemas de Steiner e as técnicas de amalgamação
utilizadas em construções do tipo partido. Esse estudo forneceu experiência
com argumentos construtivos em grafos e hipergrafos, bem como com a
organização de provas Ramsey em várias etapas.

Na segunda parte do trabalho, o foco passou à Teoria Canônica de Ramsey.
Nesse contexto, em vez de considerar apenas colorações com um número fixo
de cores, estudam-se colorações arbitrárias e os padrões estruturados que
inevitavelmente aparecem em subestruturas suficientemente grandes.

Como resultado principal da dissertação, foi obtida uma construção
canônica para hipergrafos completos ordenados. Mais precisamente, para
inteiros $n$ e $k$ com
\[
  n\geq 2k+1,
\]
foi construído um $k$-grafo ordenado $\mathcal G$, livre de
$K_{n+1}^{(k)}$, com a propriedade de que toda coloração arbitrária de
suas arestas contém uma cópia ordenada de $K_n^{(k)}$ cuja coloração
restrita é canônica.

O desenvolvimento desse resultado exigiu uma compreensão detalhada dos
padrões canônicos possíveis, do papel desempenhado pela ordenação dos
vértices e de técnicas de amalgamação de hipergrafos. Em particular, o
trabalho forneceu ao candidato experiência prévia diretamente relacionada
a duas das linhas propostas neste projeto: Teoria de Ramsey canônica e
problemas em que a estrutura da coloração possui importância tão grande
quanto a estrutura do grafo subjacente.

O projeto de doutorado representa uma continuação natural dessa formação,
mas introduz uma componente nova e essencialmente probabilística. O
objetivo agora é estudar propriedades Ramsey e anti-Ramsey em grafos
aleatórios e determinar as escalas de probabilidade em que determinados
padrões de coloração passam a ser inevitáveis.

Desde o início do doutorado, o candidato vem aprofundando sua formação em
Combinatória Probabilística e Teoria de Grafos Aleatórios, particularmente
no estudo de limiares Ramsey. Entre os resultados e técnicas estudados
estão o Teorema de Ramsey aleatório de Rödl--Ruciński, as provas e
frameworks desenvolvidos por Nenadov, Steger e colaboradores, métodos para
$0$-enunciados Ramsey, propriedades anti-Ramsey e os recentes resultados
sobre colorações canônicas em grafos aleatórios.

Esse período inicial também deu origem às primeiras investigações
diretamente relacionadas ao presente projeto. Em particular, o candidato
vem estudando a versão Ramsey canônica assimétrica, na qual os grafos-alvo
nas configurações monocromática, arco-íris e lexicográfica podem ser
distintos. Os primeiros casos considerados envolvem cliques pequenos e
têm servido tanto para identificar possíveis limiares quanto para
desenvolver construções determinísticas de coloração que poderão ser
utilizadas nos respectivos $0$-enunciados.

Paralelamente, o candidato vem estudando o problema anti-Ramsey em grafos
aleatórios e, em especial, os métodos utilizados para transformar o
$0$-enunciado probabilístico em problemas determinísticos de decomposição
e coloração. Essa investigação motivou a linha do projeto dedicada a
grafos de cintura grande.

A formação adquirida no mestrado, juntamente com o estudo desenvolvido no
período inicial do doutorado, fornece assim uma base adequada para os
problemas propostos neste projeto. Ao longo do doutorado, espera-se que o
candidato amplie essa formação, sobretudo em métodos probabilísticos,
teoria de grafos aleatórios, técnicas de decomposição e métodos de
containers, adquirindo progressivamente maior independência na formulação
e resolução de problemas de pesquisa.

% ============================================================
% CRONOGRAMA E PLANO DE TRABALHO
% ============================================================

\section{Cronograma e plano de trabalho}
\label{sec:cronograma}

O doutorado do candidato foi iniciado em julho de 2025 e tem término
previsto para julho de 2029. O projeto será desenvolvido ao longo desse
período combinando formação acadêmica, estudo da literatura, investigação
dos problemas apresentados na Seção~\ref{sec:anti-ramsey}, preparação e
submissão de artigos científicos, participação em eventos e, em etapa
posterior, redação da tese.

Embora as três linhas de pesquisa apresentadas neste projeto estejam
fortemente relacionadas, não se pretende desenvolvê-las de maneira
completamente sequencial. Problemas de diferentes linhas poderão ser
investigados simultaneamente, sobretudo quando compartilharem técnicas
ou quando resultados obtidos em uma delas sugerirem aplicações nas
demais. Ainda assim, o cronograma abaixo descreve a distribuição
principal das atividades previstas.

\subsection{Primeiro ano: julho de 2025 a julho de 2026}

O primeiro ano do doutorado foi dedicado principalmente à consolidação
da formação necessária para o desenvolvimento do projeto e ao início das
investigações em Ramsey canônico assimétrico.

Nesse período, o candidato aprofundou o estudo de Teoria de Grafos,
Combinatória Extremal e métodos probabilísticos, com particular atenção
aos resultados fundamentais sobre limiares Ramsey em grafos aleatórios.
Entre os tópicos estudados estão o Teorema de Ramsey aleatório de
Rödl--Ruciński, os frameworks para $0$-enunciados desenvolvidos por
Nenadov, Steger e colaboradores, resultados anti-Ramsey em grafos
aleatórios e o recente teorema Ramsey canônico de Kamčev--Schacht.

Paralelamente, foi iniciada a primeira frente ativa de pesquisa do
projeto: a propriedade Ramsey canônica assimétrica. O foco inicial foi
a compreensão de triplas de cliques pequenos e das ferramentas
determinísticas necessárias para construir colorações abaixo das escalas
candidatas.

Essa etapa incluiu:
\begin{itemize}
  \item estudo sistemático da literatura sobre Ramsey canônico e
        Ramsey aleatório;

  \item estudo do framework de blocos para $0$-enunciados;

  \item formulação da propriedade Ramsey canônica assimétrica;

  \item investigação inicial de triplas envolvendo
        $K_3$, $K_4$ e $K_5$;

  \item desenvolvimento de construções baseadas em degenerescência,
        templates finitos e colorações de vizinhanças;

  \item início do estudo da literatura recente sobre limiares
        anti-Ramsey.
\end{itemize}

Durante esse período, o candidato também cursou disciplinas de
pós-graduação relacionadas à sua formação matemática e computacional,
incluindo Análise de Algoritmos e Métodos Algébricos em Combinatória,
além de participar regularmente de seminários e atividades de pesquisa
em Combinatória no IME-USP.

\subsection{Segundo ano: julho de 2026 a julho de 2027}

No segundo ano, o objetivo principal será consolidar os primeiros
resultados sobre Ramsey canônico assimétrico e ampliar progressivamente
a investigação para famílias gerais de cliques.

Uma primeira meta será organizar os casos
\[
  (K_r,K_s,K_t)
\]
segundo a posição ocupada pelo maior dos três parâmetros, separando os
regimes que seguem diretamente de resultados Ramsey ou anti-Ramsey dos
casos que exigem novas construções.

Pretende-se, em particular:
\begin{itemize}
  \item finalizar a análise dos primeiros casos pequenos, incluindo
        as triplas envolvendo $K_3$, $K_4$ e $K_5$;

  \item investigar famílias gerais de triplas de cliques e buscar
        uma classificação dos regimes em que
        \[
          n^{-1/m_2(K_{\max\{r,s,t\}})}
        \]
        fornece a escala correta;

  \item desenvolver de maneira sistemática as técnicas de links,
        degenerescência e templates de coloração;

  \item estudar a possibilidade de uma redução geral
        probabilístico-determinística para a propriedade Ramsey canônica
        assimétrica;

  \item preparar para publicação os resultados obtidos nessa linha.
\end{itemize}

Em paralelo, será intensificada a investigação da linha anti-Ramsey.
A primeira meta será atacar o problema determinístico para grafos de
cintura suficientemente grande, inicialmente sob hipóteses mais fortes,
como
\[
  \operatorname{girth}(H)>8.
\]
Pretende-se estudar decomposições em florestas, subgrafos degenerados
e partes de grau limitado que permitam construir colorações próprias
sem cópias arco-íris de $H$.

Os resultados obtidos nessa etapa deverão indicar quais aspectos da
hipótese de cintura são realmente necessários e orientar tentativas
posteriores de reduzir a condição até
\[
  \operatorname{girth}(H)>4.
\]

Também será iniciado de maneira mais sistemática o estudo da terceira
linha do projeto, especialmente dos resultados recentes sobre Ramsey
canônico para ciclos.

\subsection{Terceiro ano: julho de 2027 a julho de 2028}

O terceiro ano será dedicado principalmente ao aprofundamento das duas
direções estruturais mais amplas do projeto: anti-Ramsey para grafos
gerais e Ramsey canônico além dos cliques.

Na linha anti-Ramsey, pretende-se:
\begin{itemize}
  \item aperfeiçoar as decomposições desenvolvidas no segundo ano;

  \item reduzir progressivamente a hipótese sobre a cintura do
        grafo-alvo;

  \item investigar classes intermediárias, como grafos regulares,
        bipartidos ou quase regulares de cintura grande;

  \item estudar possíveis extensões para grafos não estritamente
        $2$-balanceados;

  \item procurar e classificar eventuais obstruções determinísticas
        de baixa densidade.
\end{itemize}

Na linha Ramsey canônica para grafos gerais, o foco inicial será nos
ciclos. Pretende-se:
\begin{itemize}
  \item investigar o limiar Ramsey canônico de ciclos ímpares,
        começando pelos casos $C_5$ e $C_7$;

  \item estudar o problema para ciclos pares e a possibilidade de
        remover fatores adicionais presentes nos resultados conhecidos;

  \item compreender o papel da ordenação dos vértices;

  \item comparar diferentes ordenações de um mesmo grafo;

  \item investigar condições suficientes para que um grafo
        estritamente $2$-balanceado possua limiar canônico na escala
        $n^{-1/m_2(H)}$.
\end{itemize}

Caso os resultados obtidos nas triplas de cliques indiquem a existência
de um parâmetro de densidade assimétrico natural, essa etapa será
também utilizada para testar tal parâmetro em famílias menos
homogêneas, particularmente em problemas envolvendo cliques e ciclos.

\subsection{Estágio de pesquisa no exterior}

Planeja-se realizar, preferencialmente durante o terceiro ano do
doutorado, um estágio de pesquisa no exterior por meio de Bolsa Estágio
de Pesquisa no Exterior (BEPE).

A escolha do pesquisador anfitrião e da instituição será feita de acordo
com o desenvolvimento do projeto e com a linha que se mostrar mais
promissora ao final do segundo ano. Pretende-se buscar um grupo com
forte atuação em Teoria de Ramsey, Combinatória Probabilística, grafos
aleatórios ou métodos de containers.

O estágio deverá permitir ao candidato trabalhar de maneira intensiva
com pesquisadores especialistas nas técnicas utilizadas no projeto,
além de ampliar sua inserção internacional e possibilitar o
desenvolvimento de colaborações de longo prazo.

Um plano específico para o estágio será elaborado quando estiverem mais
claros os problemas que deverão ser atacados durante esse período.

\subsection{Quarto ano: julho de 2028 a julho de 2029}

No último ano, o candidato deverá concentrar seus esforços na
consolidação dos resultados desenvolvidos ao longo do doutorado e na
redação da tese.

As atividades previstas incluem:
\begin{itemize}
  \item conclusão dos problemas em andamento;

  \item preparação e submissão dos artigos correspondentes aos
        resultados obtidos;

  \item investigação de extensões naturais que possam ser tratadas
        com as técnicas desenvolvidas;

  \item organização dos resultados das três linhas em uma narrativa
        coerente para a tese;

  \item redação e revisão da tese de doutorado;

  \item apresentação dos resultados em seminários e eventos
        científicos;

  \item preparação para a defesa da tese.
\end{itemize}

Espera-se que, nessa etapa, as diferentes linhas de investigação
permitam identificar princípios comuns entre os problemas Ramsey
canônico, Ramsey canônico assimétrico e anti-Ramsey. Em particular,
pretende-se avaliar em que medida os resultados obtidos podem ser
organizados sob um princípio geral que relacione limiares probabilísticos
a problemas determinísticos de coloração e densidade.

Além da conclusão dos trabalhos iniciados anteriormente, o último ano
deverá permitir ao candidato explorar extensões mais ambiciosas, como
problemas para famílias de grafos, colorações localmente limitadas ou,
quando as técnicas desenvolvidas permitirem, versões para hipergrafos.

\subsection{Síntese do cronograma}

A Tabela~\ref{tab:cronograma} apresenta uma síntese das atividades
previstas ao longo do doutorado.

\begin{table}[ht]
  \centering
  \small
  \begin{tabular}{p{7.1cm}cccc}
    \hline
    \textbf{Atividade}
     & \textbf{Ano 1}
     & \textbf{Ano 2}
     & \textbf{Ano 3}
     & \textbf{Ano 4}                                     \\
    \hline

    Formação e estudo da literatura
     & $\bullet$      & $\bullet$ &           &           \\

    Ramsey canônico assimétrico: casos pequenos
     & $\bullet$      & $\bullet$ &           &           \\

    Ramsey canônico assimétrico: famílias gerais
     &                & $\bullet$ & $\bullet$ &           \\

    Anti-Ramsey e grafos de cintura grande
     & $\bullet$      & $\bullet$ & $\bullet$ &           \\

    Ramsey canônico para ciclos
     &                & $\bullet$ & $\bullet$ & $\bullet$ \\

    Ramsey canônico para grafos gerais
     &                &           & $\bullet$ & $\bullet$ \\

    Desenvolvimento de frameworks e extensões
     &                & $\bullet$ & $\bullet$ & $\bullet$ \\

    Estágio de pesquisa no exterior
     &                &           & $\bullet$ &           \\

    Preparação e submissão de artigos
     &                & $\bullet$ & $\bullet$ & $\bullet$ \\

    Participação em eventos científicos
     & $\bullet$      & $\bullet$ & $\bullet$ & $\bullet$ \\

    Redação da tese
     &                &           &           & $\bullet$ \\

    \hline
  \end{tabular}
  \caption{Cronograma geral das atividades previstas para o doutorado.}
  \label{tab:cronograma}
\end{table}


% ============================================================
% AMBIENTE DE TRABALHO E OUTROS APOIOS
% ============================================================

\section{Ambiente de trabalho e outros apoios}
\label{sec:ambiente}

O presente projeto será desenvolvido no Instituto de Matemática e
Estatística da Universidade de São Paulo, em um ambiente de pesquisa
particularmente favorável ao desenvolvimento de problemas em Teoria de
Ramsey, Combinatória Extremal e Probabilística e Teoria de Grafos
Aleatórios.

O candidato integra o grupo Comb$\Theta$ --- Combinatória, Otimização e
Teoria da Computação --- do IME-USP. O grupo reúne pesquisadores,
pós-doutorandos e estudantes que atuam em diferentes áreas da
Combinatória e mantém uma programação regular de seminários, minicursos,
workshops e visitas de pesquisadores nacionais e internacionais. Em
particular, os seminários do Comb$\Theta$ contam regularmente com
apresentações de membros do grupo e de pesquisadores visitantes,
proporcionando aos estudantes contato contínuo com problemas e técnicas
recentes da área.

O projeto será desenvolvido sob orientação de Yoshiharu Kohayakawa e
coorientação de Guilherme Oliveira Mota. Ambos possuem ampla experiência
nos temas diretamente relacionados à proposta. Kohayakawa desenvolve
pesquisa em Combinatória Extremal e Probabilística, Teoria de Ramsey,
grafos aleatórios e métodos probabilísticos, enquanto Mota atua
principalmente em Teoria de Ramsey, colorações de grafos, decomposições
e métodos probabilísticos. Os dois pesquisadores possuem uma longa
história de colaboração em problemas Ramsey e anti-Ramsey, incluindo
trabalhos sobre limiares em grafos aleatórios.

Atualmente, o ambiente de pesquisa do candidato é apoiado por diferentes
projetos que possuem forte interseção com as linhas propostas neste
plano. Destaca-se o Projeto Temático FAPESP
\emph{Combinatória clássica, assintótica, quântica e geométrica}
(Processo 2023/03167-5), coordenado por Yoshiharu Kohayakawa e vigente
de abril de 2025 a março de 2030. A vertente de Combinatória
Assintótica do projeto inclui problemas em Combinatória Extremal,
grafos aleatórios e métodos probabilísticos e reúne pesquisadores de
diversas instituições brasileiras e estrangeiras.

Entre os pesquisadores associados ao Projeto Temático estão Guilherme
Oliveira Mota, Robert Morris (IMPA), Mathias Schacht (Universität
Hamburg), Patrick Morris (Universitat Politècnica de Catalunya), Hiep
Han (Universidad de Santiago de Chile), Nicolás Sanhueza Matamala
(Universidad de Concepción), Fábio Botler (IME-USP), Carlos Hoppen
(UFRGS), Marcelo Campos (IMPA), Walner Mendonça (UFC), entre diversos
outros pesquisadores com atuação em Combinatória Extremal,
Probabilística e Teoria de Ramsey.

O Projeto Temático prevê, entre suas atividades, a realização de
workshops de pesquisa, visitas e colaborações entre os diferentes
grupos participantes. Essas atividades já vêm ocorrendo de forma
regular. Nos primeiros workshops do projeto, grupos de pesquisadores
trabalharam diretamente em problemas de Teoria de Ramsey e grafos
aleatórios.

Em particular, no segundo Workshop de Pesquisa do Projeto Temático,
realizado no IME-USP em janeiro de 2026, uma das linhas discutidas foi
uma variante assimétrica da Teoria de Ramsey canônica intimamente
relacionada ao presente projeto. O ponto de partida considerado foi o
problema no qual se buscam três grafos distintos nas configurações
monocromática, arco-íris e canônica, sendo o primeiro caso não trivial
estudado
\[
  (H_1,H_2,H_3)
  =
  (K_3,K_4,K_5).
\]
Participaram dessas discussões, entre outros, Yoshiharu Kohayakawa,
Guilherme Mota e Walner Mendonça. Parte dessas investigações foi
posteriormente continuada em um pequeno workshop realizado no IMPA,
com a participação de pesquisadores do Projeto Temático e estudantes
de doutorado do IME-USP.

Essa proximidade entre as atividades já desenvolvidas pelo grupo e uma
das linhas centrais deste projeto representa uma vantagem importante:
os problemas propostos não serão investigados isoladamente, mas dentro
de uma rede ativa de pesquisadores que já trabalha em questões
relacionadas.

Além do Projeto Temático, Guilherme Mota coordena atualmente o Auxílio
Regular FAPESP
\emph{Teoria de Ramsey: de estruturas monocromáticas a canônicas}
(Processo 2024/13859-4), vigente de maio de 2025 a abril de 2028.
Esse projeto possui relação direta com duas das principais linhas desta
proposta e reúne pesquisadores com atuação recente em diferentes
vertentes da Teoria de Ramsey, entre eles Walner Mendonça, Tássio Naia,
Patrick Morris, Marcelo Campos, Letícia Mattos e Luiz Moreira.

Mota coordena também o projeto Universal CNPq
\emph{Consolidação de avanços recentes em Teoria de Ramsey}
(Processo 420838/2025-2), vigente de novembro de 2025 a outubro de
2028. Por sua vez, Yoshiharu Kohayakawa possui bolsa de produtividade
em pesquisa do CNPq vinculada ao projeto
\emph{Combinatória assintótica de grafos, hipergrafos, estruturas
  aditivas e configurações de pontos}, vigente até 2029.

Assim, o presente projeto está inserido em diversas iniciativas de
pesquisa diretamente relacionadas aos seus objetivos científicos,
garantindo não apenas infraestrutura institucional, mas também
oportunidades frequentes de interação com pesquisadores especialistas
nos problemas e técnicas relevantes.

A rede de colaborações dos orientadores inclui alguns dos principais
pesquisadores brasileiros em Combinatória Extremal e Probabilística.
Entre eles destacam-se Robert Morris e Marcelo Campos, do IMPA; Walner
Mendonça e Fabricio Benevides, da UFC; Carlos Hoppen, da UFRGS; Fábio
Botler e Maurício Collares, além de pesquisadores de outras
instituições brasileiras. Essa interação é fortalecida pela realização
regular de eventos nacionais, como a Escola Brasileira de Combinatória,
o Workshop Brasileiro de Combinatória e o WoPOCA.

A Escola Brasileira de Combinatória constitui um exemplo particularmente
relevante desse ambiente de integração. Sua segunda edição, realizada
no IMPA em março de 2025 e organizada por Guilherme Mota, Robert Morris
e outros pesquisadores brasileiros, reuniu aproximadamente 140
participantes de nove países e de todas as regiões do Brasil, incluindo
sessões de trabalho dedicadas à discussão de problemas em aberto.

Há também forte interação com pesquisadores estrangeiros que atuam em
áreas diretamente relacionadas ao projeto. O Projeto Temático inclui,
entre outros, Mathias Schacht, Patrick Morris, Hiep Han e Nicolás
Sanhueza Matamala. O Comb$\Theta$ recebe regularmente pesquisadores
estrangeiros para visitas e seminários; nos últimos anos, por exemplo,
Julia Böttcher e Peter Allen realizaram visita ao grupo.

Outro exemplo particularmente próximo dos problemas deste projeto é
Joseph Hyde, da University of Victoria, que realizou visita de pesquisa
ao IME-USP e apresentou resultados recentes sobre limiares Ramsey
aleatórios, incluindo a Conjectura de Kohayakawa--Kreuter e problemas
anti-Ramsey e Ramsey restrito. Esses trabalhos possuem conexão direta
com os frameworks probabilístico-determinísticos que serão utilizados
nesta proposta.

A rede internacional relacionada a esses problemas inclui ainda
Natasha Morrison, atualmente na University of Victoria, pesquisadora em
Combinatória Extremal e Probabilística, grafos aleatórios e Teoria de
Ramsey. Morrison já realizou estágio de pós-doutorado no IMPA sob
supervisão de Robert Morris e é coautora de trabalhos recentes sobre
problemas Ramsey restritos e anti-Ramsey, justamente em temas próximos
à primeira linha deste projeto.

O candidato já teve oportunidade de interagir com diversos
pesquisadores dessa comunidade em seminários, escolas, workshops e
outros eventos científicos. Em particular, Walner Mendonça integrou a
comissão julgadora de sua dissertação de mestrado e mantém colaboração
científica próxima com os orientadores do presente projeto. A
continuidade dessas interações durante o doutorado deverá contribuir
significativamente tanto para o desenvolvimento dos problemas propostos
quanto para a formação científica do candidato.

Esse ambiente também é especialmente apropriado para a realização do
estágio de pesquisa no exterior previsto no cronograma. A ampla rede de
colaboradores dos orientadores oferece diversas possibilidades de
instituições e pesquisadores anfitriões com atuação diretamente
relacionada às linhas que vierem a apresentar maior desenvolvimento ao
longo do doutorado.

Consideramos, portanto, que o IME-USP e o grupo Comb$\Theta$ oferecem
um ambiente científico particularmente adequado para a execução deste
projeto. A combinação entre orientação especializada, projetos de
pesquisa vigentes, seminários regulares, workshops, visitantes e uma
ampla rede nacional e internacional de colaboradores deverá proporcionar
ao candidato as condições necessárias para desenvolver pesquisa de
fronteira nas áreas propostas.


% ============================================================
% REFERÊNCIAS
% ============================================================

\bibliographystyle{amsplain}
\bibliography{refs}

\end{document}